\subsection{A restricted dimension induction}
\label{prog:dimension-induction}

All generalized nef data in this subsection are globally nef on a fixed
compact K\"ahler carrier. For the degree arguments we use weak NQC:
a positive combination of rational degree-two classes having
nonnegative degree on compact curves. Strong NQC implies this
condition, and the contraction assertion admits this weak form
\cite[Definition~2.43]{HX}. The nef adjoint itself remains nef in
Bott--Chern cohomology throughout. A modified-big boundary means the total trace
$B+\mathbf M_X$ is globally modified big. Relative assertions retain
these absolute hypotheses. A marked model records its bimeromorphic map
from the specified source; two markings agree only under an isomorphism
commuting with those maps. Weak models below may be arbitrary normal compact K\"ahler weak log
canonical models. Their generalized nef traces are understood as
closed currents pushed forward from the fixed nef carrier; only the
whole adjoint is required to be a Bott--Chern class. In particular no
factoriality, and no separate Bott--Chern class for either $K_Y+B_Y$
or the nef trace, is required of an auxiliary weak target. This
category contains both the working models and the adjunction strata
used in special termination. Chosen program outputs remain globally
strongly $\mathbb Q$-factorial.

For dimension at most $d$, write $\mathsf B_d$ for semiampleness of a
nef gklt adjoint with modified-big boundary, including a Moishezon
connected contraction onto a normal compact K\"ahler target. It makes
no assertion that every such contraction is projective.
Write $\mathsf C_d$ for the NQC non-klt contraction assertion of
\cite[Theorem~1.6]{HX}, and $\mathsf C_{d,\mathrm{big}}$ for its
big-adjoint case. The relative assertions are:
\begin{itemize}
\item $\mathsf M_d$: a relatively pseudo-effective gklt adjoint over
any proper map to a normal compact K\"ahler base $S$, with the globally
nef and modified-big data just specified, has a chosen nonextracting
good log terminal model over $S$. Its working space is globally
strongly $\mathbb Q$-factorial and compact K\"ahler; its connected
relative canonical morphism has normal compact K\"ahler target and
is Moishezon. From a smooth common carrier,
all Bott--Chern classes have forward traces, and the chosen model has
the rational degree-two and Bott--Chern exceptional splittings proved
below.
\item $\mathsf F_d$: a compact polytope of these data on one fixed
carrier has a polyhedral relative effective locus, finitely many
relative canonical chambers with chosen good terminal models, and
finitely many marked weak models, accounting for every weak model
in the category specified above.
\end{itemize}
Passing to the proper image and then to its normal Stein factor allows
all base maps to be taken surjective with connected fibres. These
operations preserve the compact K\"ahler base category.

The independent inputs are the analytic cone theorem, the projective
results of Proposition~\ref{prog:projective-good-model}, relative
vanishing, adjunction, relative-canonical positivity, and the
nef-and-big K\"ahler criterion. The base presentations and
multiplier-ideal descent used in the induction are proved below.
The projective constructions use the actual-line lemmas of the preceding
subsection. No assertion of termination of an arbitrary sequence of
flips is used.

\subsubsection{Base presentations and multiplier ideals}

The following constructions take place in the category of generalized
data specified above. In particular, the whole adjoint is a
Bott--Chern class, whereas the boundary and nef traces on a singular
space are interpreted as currents. On a smooth carrier the structure
boundary is a real divisor; its negative components may be exceptional
over the space of the pair. The nef part is specified on that carrier
and pulled back to higher models.

\begin{lemma}[Adjoint-preserving preparation]
\label{prog:base-presentation-preparation}
Let $(X,B+\mathbf M)$ be a generalized pair on a normal compact
K\"ahler space, with globally modified-big total boundary
$B+\mathbf M_X$. Write
\[
 A=[K_X+B+\mathbf M_X],\qquad
 \mathcal J_o=\mathcal J(X,B+\mathbf M).
\]
There are generalized data on $X$ with the same adjoint class, a
projective carrier $p:U\to X$ with $U$ smooth and compact K\"ahler, an SNC real
divisor $C$ whose negative components are $p$-exceptional, and a
K\"ahler class $\beta$ such that
\begin{equation}\label{prog:base-presentation-prepared}
 p^*A=[K_U+C]+\beta,\qquad
 \mathcal J_1=p_*\mathcal O_U(-\lfloor C\rfloor)
                 \supseteq\mathcal J_o.
\end{equation}
The carrier may be required to dominate any prescribed projective
modification of $X$, with the structure and exceptional divisors
having simultaneous SNC support. The prepared total boundary is
globally modified big. In particular, the preparation preserves the
gklt condition.
\end{lemma}

\begin{proof}
Choose a common projective log resolution $p_0:U_0\to X$ of the
original nef carrier and the prescribed modification, and write
\[
 p_0^*A=[K_{U_0}+C^0]+\mathbf M_{U_0}.
\]
The class $\mathbf M_{U_0}$ is nef, and the negative part $(C^0)^-$
is exceptional over $X$. The modified-big preparation
\cite[Lemmas~2.10 and~2.23]{HX} can be made with the prescribed
resolution. We give its divisor and ideal details. Choose an effective divisor $E_0$ with sufficiently large coefficients
and full $p_0$-exceptional support so that the class
\[
 [(C^0)^+]+\mathbf M_{U_0}+[E_0]
\]
is big. Resolve a Kodaira decomposition of this class. On a further
projective modification $a:U\to U_0$ with $U$ smooth one obtains
\[
 a^*\bigl([(C^0)^+]+\mathbf M_{U_0}+[E_0]\bigr)
       =\kappa+[D],\qquad D\geq0,
\]
where $\kappa$ is K\"ahler. The support can be made SNC together
with all the other divisors. On any additional blowup, subtract a
sufficiently small effective exceptional class from the pulled-back
K\"ahler class and add it to the divisor; an exceptional divisor
whose negative is relatively ample makes the resulting class
K\"ahler again.

Put $p=p_0a$ and
\[
 C^{\mathrm{alt}}=D-a^*(C^0)^--a^*E_0-K_{U/U_0}.
\]
Then
\[
 p^*A=[K_U+C^{\mathrm{alt}}]+\kappa,
\]
and every negative component of $C^{\mathrm{alt}}$ is
$p$-exceptional. Let $C^{\mathrm{old}}$ be the crepant structure
boundary of the original data on this same model. For $0<t\ll1$,
set
\[
 C=(1-t)C^{\mathrm{old}}+tC^{\mathrm{alt}},\qquad
 \beta=(1-t)a^*\mathbf M_{U_0}+t\kappa.
\]
The class $\beta$ is K\"ahler, the adjoint identity is unchanged,
and negative coefficients still occur only on exceptional divisors.
There are only finitely many coefficients. Thus $t$ can be chosen so
that
\[
 \lfloor C\rfloor\leq\lfloor C^{\mathrm{old}}\rfloor.
\]
Pushing forward the corresponding inclusion of invertible sheaves
gives~\eqref{prog:base-presentation-prepared}. The K\"ahler class
$\beta$ plus the strict transform of the effective trace boundary is
big; its pushforward is the prepared total boundary. If the original pair
is gklt, its multiplier ideal is the unit ideal, and hence so is
$\mathcal J_1$.
\end{proof}

\begin{proposition}[Base presentations for projective contractions]
\label{prog:base-presentation}
Let $f:X\to Z$ be a projective contraction between normal irreducible
compact K\"ahler spaces. Suppose that $(X,B+\mathbf M)$ has globally
modified-big total boundary and that
\[
 A=[K_X+B+\mathbf M_X]=f^*\gamma,
 \qquad \gamma\in H^{1,1}_{\mathrm{BC}}(Z).
\]
There is an adjoint-preserving preparation as in
Lemma~\ref{prog:base-presentation-preparation}, with multiplier ideal
$\mathcal J_1\supseteq\mathcal J_o$, for which
\begin{equation}\label{prog:base-presentation-Stein}
 R^if_*\mathcal J_1=0\quad(i>0),\qquad
 f_*\mathcal O_{N_1}=\mathcal O_Z/(f_*\mathcal J_1),
 \qquad N_1=V(\mathcal J_1).
\end{equation}
If $N_1$ dominates $Z$, then $f_*\mathcal J_1=0$.
Otherwise there are generalized data $(Z,B_Z+\mathbf N)$, a
projective carrier $\tau:T\to Z$ with $T$ smooth and compact K\"ahler, a real
SNC structure divisor $C_T$, and a K\"ahler class $\beta_T$ such that
\begin{equation}\label{prog:base-presentation-identities}
 \begin{gathered}
 [K_Z+B_Z+\mathbf N_Z]=\gamma,\qquad
 \tau^*\gamma=[K_T+C_T]+\beta_T,\qquad
 \mathbf N_T=\beta_T,\\
 B_Z=\tau_*C_T\geq0,\qquad
 B_Z+\mathbf N_Z\ \text{is globally modified big}.
 \end{gathered}
\end{equation}
The negative components of $C_T$ are $\tau$-exceptional, and
\begin{equation}\label{prog:base-presentation-ideal}
 \mathcal J_Z:=\mathcal J(Z,B_Z+\mathbf N)
       =\tau_*\mathcal O_T(-\lfloor C_T\rfloor)
       \supseteq f_*\mathcal J_1.
\end{equation}
In particular the base data are gklt wherever $f_*\mathcal J_1$ is
the unit ideal. If the source data are gklt, the base data are gklt.
All adjoint identities are preserved by pullback to higher models.
\end{proposition}

\begin{proof}
Choose a projective resolution $\mu:W\to Z$, with $W$ smooth
compact K\"ahler and reduced exceptional divisor $F$ having SNC
support. Apply Lemma~\ref{prog:base-presentation-preparation} on a
carrier dominating the main component of the base change. We obtain
projective maps $p:U\to X$ and $q:U\to W$ with $fp=\mu q$ and
\begin{equation}\label{prog:base-presentation-source}
 [K_U+C]+\beta=p^*A=q^*\mu^*\gamma,
 \qquad \mathcal J_1=p_*\mathcal O_U(-\lfloor C\rfloor).
\end{equation}
Here $U$ is smooth compact K\"ahler, $\beta$ is K\"ahler, and all
negative components of $C$ are $p$-exceptional. The maps are
projective by composition and factorization over the separated base
$W$. The map $q$ has connected fibres: this holds over the open set
where $\mu$ is an isomorphism, and its Stein factorization is finite
and bimeromorphic over the normal space $W$.

The integral divisor $-\lfloor C\rfloor$ has class
\[
 [-\lfloor C\rfloor]
   =[K_U+\{C\}]+\beta-(fp)^*\gamma.
\]
The pair $(U,\{C\})$ is klt with SNC boundary. Relative
Kawamata--Viehweg vanishing, first for $p$ and then for $fp$, applies
because the positive class modulo pullbacks is K\"ahler
\cite[Theorem~2.2]{HX}. Equivalently, the real line bundle
$-K_U-C$ is relatively ample,
so the vanishing follows from \cite[Theorem~1.2]{FujinoVanishing2025}
with the identity as its first morphism. Leray gives $R^if_*\mathcal J_1=0$ for
$i>0$. Applying $f_*$ to the sequence defining $N_1$, and using
$f_*\mathcal O_X=\mathcal O_Z$, proves
\eqref{prog:base-presentation-Stein}. If $N_1$ dominates $Z$, the
ideal $f_*\mathcal J_1\subseteq\mathcal O_Z$ vanishes generically
and therefore vanishes. This proves that alternative.

Suppose from now on that $N_1$ does not dominate $Z$. Every
$q$-horizontal coefficient of $C$ is then less than one. Indeed,
a component with coefficient at least one forces the multiplier
ideal to be proper along its image in $X$, and its image would make
$N_1$ dominate $Z$. If $Z$ is a point, the asserted base presentation
is the dimension-zero presentation. We therefore assume
$b=\dim Z>0$.

\paragraph{The fibre metric.}
Choose a smooth representative $\rho$ of
$\mu^*\gamma-c_1(K_W)$, a K\"ahler form $\omega_W$, and a
K\"ahler form representing $\beta$, still denoted by $\beta$.
For some $a>0$,
\[
 \beta\geq a q^*\omega_W.
\]
Consider the holomorphic line bundle
\[
 L=\mathcal O_U(-K_{U/W}-\lfloor C\rfloor).
\]
Equation~\eqref{prog:base-presentation-source} gives a singular
Hermitian metric $h$ on $L$ whose curvature current is
\begin{equation}\label{prog:base-presentation-curvature}
 \Theta_h(L)=\beta+[\{C\}]-q^*\rho.
\end{equation}
Its weights are smooth functions plus the divisorial logarithms
with coefficients in $\{C\}$. We normalize $dd^c$ so that
$dd^c\log|z|^2=[z=0]$ and write squared-norm weights as $e^{-\varphi}$.
Let $e$ be the canonical meromorphic section of
\[
 K_{U/W}\otimes L=\mathcal O_U(-\lfloor C\rfloor),
\]
so that $\operatorname{div}(e)=-\lfloor C\rfloor$.

Choose an analytic Zariski open set $W^\circ\subseteq W$ on which
$q$ is a submersion, the horizontal boundary is relatively SNC,
there are no vertical boundary components, $\mu$ is an isomorphism
onto an open subset of $Z$, and coherent base change holds in degree
zero. On $q^{-1}(W^\circ)$, the divisor $-\lfloor C\rfloor$ is
effective and exceptional over $X$. Normality and the contraction
property therefore identify
\[
 q_*(K_{U/W}\otimes L)|_{W^\circ}
       =\mathcal O_{W^\circ}e.
\]
Here $K_{U/W}$ is viewed as the relative canonical line bundle. The
fibre norm
\[
 I(w)=\int_{U_w}|e_w|_h^2
\]
is positive and finite. The fractional SNC weights are integrable;
any zeros of $e_w$ only improve integrability. Relative SNC
coordinates and a partition of unity show that $I$ is smooth:
parameter derivatives act on smooth factors multiplying fixed
integrable monomial weights.

Set $\psi=-\log I$ on $W^\circ$. Then
\begin{equation}\label{prog:base-presentation-open-positivity}
 \rho+dd^c\psi\geq a\omega_W.
\end{equation}
Here are the details for the singular weights. Locally choose
$\ell$ with $dd^c\ell=\rho-a\omega_W$, and put
$h'=he^{-q^*\ell}$. Its weights are psh, and its fibre norm satisfies
$I'=e^{-\ell}I$. Restrict the family to a small coordinate line disk
$\Delta_R$, with centre $0$. The sharp $L^2$ extension theorem gives
a twisted top form extending $e_0$ and satisfying
\begin{equation}\label{prog:base-presentation-sharp-extension}
 \int_{q^{-1}(\Delta_R)}|E|_{h'}^2
        \leq\pi R^2 I'(0).
\end{equation}
Top-form norms here are normalized to use Euclidean area in the
base coordinate. We use the Stein hypersurface form of the optimal
extension theorem \cite[Theorem~2.2 and \S3.5, (3.1)]{GuanZhouOptimal2015},
with defining function the disk coordinate divided by $R$.

To check applicability, embed the restricted projective family into
projective space over the disk and remove a relative hyperplane
section that does not contain the central fibre. The complement is
Stein. A holomorphic section of $L$ on this Stein space can be chosen
nonzero at a point of the central fibre; deleting its zero divisor
trivializes $L$ and leaves a Stein space. The central fibre is smooth
and connected, hence irreducible, so both deletions leave a dense
open subset of it. For a zero-dimensional fibre choose the deletions
to avoid its point. In the resulting trivialization the weight is
psh, and the sharp extension theorem applies to canonical forms.
The extension extends across the deleted hypersurfaces in reverse
order by local $L^2$ removability: in every original frame the metric
is bounded below by a smooth positive metric. The extension retains
its norm bound and its prescribed restriction to the central fibre.

After dividing $E$ by the base differential, the rank-one direct
image identifies it with $b(t)e$, where $b$ is holomorphic and
$b(0)=1$. Thus the disk average of $|b(t)|^2I'(t)$ is at most $I'(0)$.
Jensen's inequality and the submean inequality for $\log|b|^2$ give
\[
 \log I'(0)
 \geq\operatorname{avg}_{\Delta_R}
             \bigl(\log|b(t)|^2+\log I'(t)\bigr)
 \geq\operatorname{avg}_{\Delta_R}\log I'(t).
\]
The same argument applies to every sufficiently small disk at every
centre and in every complex line direction. Hence $-\log I'=\psi+\ell$ is psh, which proves
\eqref{prog:base-presentation-open-positivity}.

\paragraph{Extension over the base.}
Let $D$ be a prime divisor in $W\setminus W^\circ$. Choose a component
$Q$ of $q^*D$ dominating $D$. When $D$ is not $\mu$-exceptional, choose
$Q$ not $p$-exceptional: take a component over the corresponding prime
on $Z$ already on $X$ and then its strict transform. Put
\[
 u=\operatorname{ord}_Q(q^*D),\qquad c=\operatorname{coeff}_Q(C).
\]
At a general point of $Q$, avoiding all other relevant divisors,
there are local coordinates with
\[
 t_1=z_1^u,\qquad t_j=z_j\quad(2\leq j\leq b).
\]
Indeed the map $Q\to D$ is submersive at a general smooth point,
and a local unit in the first equation is absorbed by a holomorphic
$u$th root. Such points can be chosen above every point of a dense
analytic open subset of $D$ by generic smoothness and properness.
Integration over a fixed small polydisk in the remaining fibre
coordinates gives
\begin{equation}\label{prog:base-presentation-divisor-bound}
 I(t)\geq c_*|t_1|^{-2(c+u-1)/u}
\end{equation}
for some $c_*>0$ near each such base point. The section and metric
contribute the exponent $c$, while division by the base differential
in the relative canonical form contributes $u-1$.

For a nonexceptional $D$ the chosen $Q$ has $c\geq0$, so
\eqref{prog:base-presentation-divisor-bound} bounds $\psi$ from above
near its general points. For the finitely many exceptional divisors,
the same estimate supplies a nonnegative integer $m$ such that, if
$k$ is a smooth metric on $\mathcal O_W(F)$ and $s_F$ its divisor
section,
\[
 \varphi=\psi+m\log|s_F|_k^2
\]
is locally bounded above near general points of every divisor.
Adding local potentials of $\rho+m c_1(k)-a\omega_W$ to
$\varphi$ gives psh weights by
\eqref{prog:base-presentation-open-positivity}. These weights
extend across the divisors and then across the remaining analytic
set of codimension at least two, giving
\begin{equation}\label{prog:base-presentation-positive-current}
 S=\rho+m c_1(k)+dd^c\varphi\geq a\omega_W
\end{equation}
on all of $W$.

\paragraph{A K\"ahler carrier.}
Regularize the current in
\eqref{prog:base-presentation-positive-current} in its class, retaining
a lower bound for the potentials. More precisely, choose an
approximation with analytic singularities such that
\begin{equation}\label{prog:base-presentation-regularization}
 S'=\rho+m c_1(k)+dd^c\varphi',\qquad
 S'\geq\frac a2\omega_W,\qquad
 \varphi'\geq\varphi-O(1).
\end{equation}
These simultaneous properties are the Bergman regularization
properties of \cite{DemaillyRegularization1992}; see also
\cite[pp.~59--60]{PopoviciRegularization2004}.
Resolve the analytic singularities and the support of $F$ by a
sequence of blowups with smooth centres $r:T\to W$, with simultaneous
SNC support, and write
\[
 r^*S'=[H]+\zeta,\qquad H\geq0.
\]
The residual current $\zeta$ has locally bounded potentials. Removing
the divisorial part from the positive current
$r^*S'-(a/2)r^*\omega_W$ shows that
\[
 \zeta\geq\frac a2r^*\omega_W.
\]
The difference has vanishing Lelong numbers, so its class is nef by
regularization. Choose an effective $r$-exceptional divisor $G$ with
$-G$ relatively ample. For $0<\eta\ll1$, the class
$(a/2)r^*\omega_W-\eta[G]$ is K\"ahler. Therefore
\begin{equation}\label{prog:base-presentation-carrier-data}
 \begin{split}
 \tau&=\mu r,\qquad
 \beta_T=[\zeta]-\eta[G],\\
 C_T&=H+\eta G-mr^*F-K_{T/W}
 \end{split}
\end{equation}
has K\"ahler $\beta_T$. If $r$ is an isomorphism, omit $G$ and its
correction. The choice of $\eta$ will be decreased once more below.

The exact class calculation is
\[
 \begin{split}
 [K_T+C_T]+\beta_T
 &=r^*c_1(K_W)+[H]+[\zeta]-m[r^*F]\\
 &=r^*(c_1(K_W)+[\rho])=\tau^*\gamma.
 \end{split}
\]
All possibly negative terms in $C_T$ are $\tau$-exceptional.
Consequently $B_Z=\tau_*C_T$ is effective. Specify the nef b-data by
$\mathbf N_T=\beta_T$; the displayed identity defines generalized
data on $Z$ with adjoint class $\gamma$. At the level of currents,
pushing the identity gives the required trace identity on $Z$:
exceptional divisor currents push to zero, and a $dd^c$ difference
pushes to a $dd^c$ difference. Canonical traces can be computed using
local meromorphic canonical frames; a change of frame changes the
corresponding Chern potentials by the usual logarithmic transition
term. Thus only the whole adjoint, rather
than its individual traces, needs to be a Bott--Chern class. The
class $\beta_T$ plus the strict transform of $B_Z$ is big on $T$,
so its pushforward $B_Z+\mathbf N_Z$ is globally modified big.
These are~\eqref{prog:base-presentation-identities}.

\paragraph{The multiplier ideal.}
Take a local holomorphic function $v\in f_*\mathcal J_1$. By the
pushforward description of $\mathcal J_1$,
\begin{equation}\label{prog:base-presentation-source-orders}
 \operatorname{div}(q^*\mu^*v)\geq\lfloor C\rfloor.
\end{equation}
Thus $q^*\mu^*v\,e$, tensored with a local base canonical frame,
is a holomorphic $L$-valued top form on $U$. Its squared norm is
locally integrable by the fractional SNC weights in
\eqref{prog:base-presentation-curvature}. Over the submersion locus,
fibre integration followed by base integration is precisely this
top-form integral. Properness and a finite coordinate covering
therefore give
\begin{equation}\label{prog:base-presentation-base-integral}
 \int |\mu^*v|^2e^{-\psi}<\infty
\end{equation}
locally on $W$; the integral is taken where the original fibre
metric is defined.

The lower bound on $\varphi'$ in
\eqref{prog:base-presentation-regularization} preserves this
integrability after replacing $\psi$ by
$\varphi'-m\log|s_F|_k^2$. Change variables through $r$. Along a
prime divisor of $T$, let $h,j,k_D,d$ be respectively the coefficients
of $H$, $\operatorname{div}(\tau^*v)$, $K_{T/W}$, and $r^*F$.
The local integrand has exponent
$2(j-h+md+k_D)$, up to a bounded positive factor. Its integrability
implies $j-h+md+k_D>-1$, or, since $j$ is integral,
\[
 j\geq\lfloor h-md-k_D\rfloor.
\]
Therefore
\[
 \operatorname{div}(\tau^*v)
       \geq\lfloor H-mr^*F-K_{T/W}\rfloor.
\]
Choose $\eta>0$ sufficiently small that adding $\eta G$ does not
change any of these finitely many round-downs. The resulting
inequality is exactly
$\operatorname{div}(\tau^*v)\geq\lfloor C_T\rfloor$, proving
\eqref{prog:base-presentation-ideal} as an inclusion of ideal sheaves.
In particular, if $f_*\mathcal J_1$ is the unit ideal on an open set,
then so is $\mathcal J_Z$, and the base data are gklt there.

Finally, the nef data descend to $T$, and the structure boundaries
on higher models are defined by crepant pullback of
\eqref{prog:base-presentation-identities}. On every common refinement
of the source and base carriers, both adjoints are the pullback of
$\gamma$. Thus the identities remain exact on higher models.
\end{proof}

\begin{lemma}[Contractions of closed subspaces]
\label{prog:base-emc}
Let $D$ be a compact complex space, possibly nonreduced, and let
$\delta\in H^2(D,\R)$ have nonnegative degree on compact curves.
Use the definitions of a Moishezon morphism and of a contraction
defined by a class in \cite[Definitions~2.32 and~2.34]{HX}.
Then the following hold.
\begin{enumerate}
\item If $\delta$ is endowed with a Moishezon contraction, so is
its restriction to every closed analytic subspace of $D$.
\item Let $g:D\to V$ be a projective contraction of compact
complex spaces, so $g_*\OO_D=\OO_V$, and suppose
$\delta=g^*\epsilon$ for a class $\epsilon\in H^2(V,\R)$
having nonnegative degree on compact curves. Then $\delta$ is
endowed with a Moishezon contraction if and only if $\epsilon$
is endowed with one.
\end{enumerate}
\end{lemma}

\begin{proof}
Write $h:D\to S$ for a Moishezon contraction defined by
$\delta$. For a closed subspace $D_0\subset D$, take the Stein
factorization
\[
                 D_0\longrightarrow S_0\longrightarrow S
\]
of its restriction. The first map is a contraction, and the second
is finite. Thus a curve in $D_0$ is contracted by the first map
exactly when its $\delta$-degree is zero. Choose a projective
surjection $P\to D$ witnessing that $h$ is Moishezon, with
$P\to S$ projective. The closed subspace
$P_0=P\times_D D_0$ is projective over $S$ and projective
and surjective over $D_0$. It is also projective over $S_0$:
a relatively ample line for $P_0\to S$ remains ample on the
fibres of $P_0\to S_0$. This proves the first assertion,
including its scheme structure.

For the descent direction of the second assertion, every curve
in a fibre of $g$ has $\delta$-degree zero. A connected
projective complex space is connected by chains of compact curves
after passing to its reduction. Therefore $h$ is constant on
the underlying set of each fibre of $g$. It factors as
$h=k\circ g$ for a holomorphic map $k:V\to S$.
To see that this factorization holds also for nonreduced spaces,
choose a coordinate neighbourhood in $S$ of the image of a fibre.
Properness permits restriction to a neighbourhood of its point
in $V$ whose inverse image maps into that coordinate
neighbourhood. The coordinate functions of $h$ then descend
uniquely through $g_*\OO_D=\OO_V$; their analytic relations
descend through the same equality. These local factorizations
glue. Furthermore,
\[
                 k_*\OO_V=h_*\OO_D=\OO_S.
\]
Every irreducible compact curve $C\subset V$ is dominated by
a compact curve $C'\subset D$. Indeed, after base change to
the normalization of $C$, take a component dominating that curve
and a projective multisection. The projection formula gives
\[
       \delta\cdot C'
          =\deg(C'\to C)\,\epsilon\cdot C.
\]
Consequently $k$ contracts exactly the curves of
$\epsilon$-degree zero. The composite $P\to D\to V$
is a projective surjection, and $P\to S$ is projective, so
it witnesses that $k$ is Moishezon.

Conversely, suppose $k:V\to S$ is a Moishezon contraction
defined by $\epsilon$. Then $k\circ g$ is a contraction,
and the projection formula verifies its curve criterion, treating
separately curves contracted by $g$. If $Q\to V$ is a
projective surjection witnessing that $k$ is Moishezon, then
$D\times_V Q\to D$ is projective and surjective, while
$D\times_V Q\to Q\to S$ is projective. This supplies the
required witness.
\end{proof}

\begin{lemma}[Crepant transport of the multiplier closed subspace]
\label{prog:multiplier-transport}
Let $g:X\to Y$ be a projective bimeromorphic morphism of
normal compact K\"ahler spaces. Suppose effective generalized
pair data on $X$ and $Y$ have the same globally nef b-part
and are crepant over $Y$. In particular their adjoint classes
$A_X$ and $A_Y$ satisfy $A_X=g^*A_Y$ in Bott--Chern
cohomology. Write $J_X,J_Y$ for their multiplier ideals and
$N_X=V(J_X)$, $N_Y=V(J_Y)$ for their non-klt closed
subspaces. Then
\[
       g_*J_X=J_Y,\qquad R^ig_*J_X=0\quad(i>0),
       \qquad g_*\OO_{N_X}=\OO_{N_Y}.
\]
The restriction $N_X\to N_Y$ is a projective contraction.
Hence, whenever $A_Y$ is nef, the restriction of $A_X$ to
$N_X$ is endowed with a Moishezon contraction if and only if
the restriction of $A_Y$ to $N_Y$ is endowed with one.
\end{lemma}

\begin{proof}
Take a projective common log resolution $r:U\to X$ carrying
the nef datum, and put $s=g\circ r$. By crepancy there is
one real SNC divisor $C$ such that
\[
            K_U+[C]+\beta=s^*A_Y=r^*A_X,
\]
where $\beta$ is globally nef. The multiplier ideals are
\[
 J_X=r_*\OO_U(-\lfloor C\rfloor),\qquad
 J_Y=s_*\OO_U(-\lfloor C\rfloor).
\]
Their direct-image equality follows at once. For vanishing, set
$L=\OO_U(-\lfloor C\rfloor)$ and
$\Delta=C-\lfloor C\rfloor$. The pair $(U,\Delta)$
is klt, and the actual real line
\[
                  L-(K_U+\Delta)=-K_U-C
\]
has class $\beta-s^*A_Y$. It is relatively nef over $Y$
and relatively big because $s$ is bimeromorphic. The same
statements hold over $X$. Relative Kawamata--Viehweg vanishing,
in the form of \cite[Theorem~1.2]{FujinoVanishing2025} with
the first morphism the identity and the real line
$\mathcal H=-K_U-C$, therefore gives
\[
 R^is_*L=R^ir_*L=0\qquad(i>0).
\]
The real-line identity in that theorem is the displayed definition
of $\mathcal H$; its relative positivity is checked by the class
identity above. The klt condition means that there are no proper log
canonical strata to test for relative log bigness. Leray gives
$R^ig_*J_X=0$ for $i>0$. Pushing forward
$0\to J_X\to\OO_X\to\OO_{N_X}\to0$ and using
$g_*\OO_X=\OO_Y$ now proves the asserted structure-sheaf
equality. It also identifies the scheme-theoretic image with
$N_Y$, giving the projective contraction between these closed
subspaces. Apply Lemma~\ref{prog:base-emc} to conclude.
\end{proof}

\begin{remark}
The preceding lemma also transports the contraction hypothesis
through a crepant projective flip: apply it to the two contractions
to their common base, equipped with the descended adjoint and
the common nef b-data. On that base the boundary is the effective
pushforward of the source boundary, and the nef data remain on
the same carrier; the descended Bott--Chern adjoint and its
identity on a common resolution define the crepant base data
in the whole-adjoint and current-trace category used here.
All assertions concern the fixed adjoint
and its multiplier ideal. No conclusion about the multiplier
ideal of a different driving adjoint is used.
\end{remark}

\begin{corollary}[The contraction step over a lower-dimensional base]
\label{prog:base-emc-induction}
Use the hypotheses and notation of
Proposition~\ref{prog:base-presentation}. Suppose that
$A=f^*\gamma$ is nef and that its restriction to the original
non-klt closed subspace $N_o=V(J_o)$ is endowed with a
Moishezon contraction. Put
\[
                I=f_*J_1,\qquad N_*=V(I)\subset Z.
\]
Then $\gamma|_{N_*}$ is endowed with a Moishezon contraction.
If the prepared non-klt closed subspace $N_1=V(J_1)$ dominates
$Z$, then $\gamma$ and $A$ are endowed with Moishezon
contractions. Otherwise the generalized pair supplied on $Z$
has a non-klt closed subspace endowed with the contraction
defined by the restricted class of $\gamma$. In this latter
case, if $\gamma$ is nef and NQC and the assertion
$\mathsf C_{\dim Z}$ is available, then $\gamma$ and $A$
are endowed with Moishezon contractions.
\end{corollary}

\begin{proof}
The inclusion $J_o\subset J_1$ gives $N_1\subset N_o$
as closed analytic subspaces. By
Lemma~\ref{prog:base-emc}, the given contraction restricts to
$N_1$. Proposition~\ref{prog:base-presentation} identifies
\[
              f_*\OO_{N_1}=\OO_Z/I=\OO_{N_*}.
\]
Thus the restriction $N_1\to N_*$ is a projective
contraction. The restricted adjoint is the pullback of
$\gamma|_{N_*}$, so Lemma~\ref{prog:base-emc} descends
the contraction to $N_*$. Nonnegative degrees on curves of
$N_*$ follow from projective multisections and nefness of
$A$.

If $N_1$ dominates the normal irreducible space $Z$, then
$I=0$: a local function in $I$ pulls back to zero on the
dominating closed subspace and hence vanishes generically on
$Z$. Normality, and in particular reducedness, gives its
vanishing. Therefore $N_*=Z$, which proves the assertion in
this case.

If $N_1$ does not dominate $Z$, the constructed base ideal
satisfies $I\subset J_Z$. Hence
\[
                 V(J_Z)\subset V(I)=N_*
\]
as closed analytic subspaces. Restriction by
Lemma~\ref{prog:base-emc} provides the required contraction
on $V(J_Z)$, with its complete scheme structure. The
generalized pair on $Z$ has globally nef carrier data and
globally modified-big total boundary by
Proposition~\ref{prog:base-presentation}. When the stated
nef and NQC conditions hold, these facts are precisely the
hypotheses of $\mathsf C_{\dim Z}$, which gives the
contraction on $Z$. Finally Lemma~\ref{prog:base-emc}
pulls it back through the projective contraction $f$ to $X$.
The case of an empty non-klt closed subspace has the same
interpretation, with its contraction condition vacuous.
\end{proof}


\subsubsection{A fixed integral grid and cohomology transport}

\begin{lemma}[Integral degree-two descent for a locally log-Fano map]
\label{prog:integral-h2}
Let $f:X\to Z$ be a projective surjective morphism with connected fibres
between normal complex analytic spaces. Assume that $f_*\mathcal O_X=
\mathcal O_Z$, that $R^if_*\mathcal O_X=0$ for $i>0$, and that every point
of $Z$ has a neighbourhood $U$ on which there is a klt pair
$(X_U,\Delta_U)$ with $-(K_{X_U}+\Delta_U)$ ample over $U$.
Then $f^*:H^2(Z,\mathbb Z)\to H^2(X,\mathbb Z)$ is injective, and its
image consists exactly of the integral classes having degree zero on
every compact curve contracted by $f$.
\end{lemma}

\begin{proof}
Properness, normality and connected fibres identify both
$f_*\mathcal O_X=\mathcal O_Z$ and
$f_*\mathcal O_X^*=\mathcal O_Z^*$. Apply $Rf_*$ to the analytic
exponential sequence. The local surjectivity of
$\mathcal O_Z\to\mathcal O_Z^*$ and the stated coherent vanishing give
\[
 R^1f_*\mathbb Z_X=0,
 \qquad
 R^1f_*\mathcal O_X^*\simeq R^2f_*\mathbb Z_X.
\]
Let $\xi\in H^2(X,\mathbb Z)$ have zero degrees on contracted curves.
A germ of its edge image in $R^2f_*\mathbb Z_X$ is, under this
isomorphism, represented after shrinking $U$ by an actual line bundle
$L$ on $X_U$. Equality of these germs means that, after further
shrinking, $c_1(L)$ and $\xi|_{X_U}$ have the same class. In particular
$L$ has degree zero on every curve in every fibre above this smaller
neighbourhood.

On a sufficiently small Stein base neighbourhood, twisting by a
relatively ample line and taking a quotient of two nonzero relative
sections gives a meromorphic Cartier representative for $L$. Thus
the Cartier-divisor form of base point freeness applies.
The line bundle $L$ is relatively nef, and
$aL-(K_{X_U}+\Delta_U)$ is relatively ample for every real $a$.
The projective analytic base-point-free theorem applies to the Cartier
line $L$. Its conclusion is generation for \emph{every} sufficiently
large integer power, not only for a divisible subsequence
\cite[Theorems~6.2 and~6.5]{FujinoAnalyticBCHM}. Thus, over a further relatively compact
neighbourhood, both $L^m$ and $L^{m+1}$ are relatively generated.
A generated line of degree zero on every curve of a projective connected
fibre defines a constant projective map on that fibre: the restriction
of the generated line is the pullback of $\mathcal O(1)$, and a
positive-dimensional projective image contains a curve of positive
degree. The corresponding relative image is therefore finite and
bimeromorphic over the normal base, and is the base itself. Consequently
\[
 L^m\simeq f^*M_m,\qquad L^{m+1}\simeq f^*M_{m+1}
\]
for line bundles on the neighbourhood in $Z$. Taking their quotient
proves $L\simeq f^*(M_{m+1}\otimes M_m^{-1})$. Its germ in
$R^1f_*\mathcal O_X^*$ is zero, so the edge image of $\xi$ is zero.

The integral Leray spectral sequence, using
$R^1f_*\mathbb Z_X=0$, gives
\[
 0\longrightarrow H^2(Z,\mathbb Z)
 \xrightarrow{f^*}H^2(X,\mathbb Z)
 \longrightarrow H^0(Z,R^2f_*\mathbb Z_X).
\]
Exactness at the middle term proves the assertion. The argument is
integral throughout and does not discard torsion. The converse follows
by restriction of a pulled-back class to a fibre.
\end{proof}

\begin{lemma}[A fixed integral grid along existing crepant steps]
\label{prog:fixed-integral-grid}
Let $X_0$ be a compact K\"ahler space, and suppose that
\[
 \alpha_0=\sum_{j=1}^r a_j\gamma_{0j},\qquad a_j>0,
 \qquad \gamma_{0j}\in H^2(X_0,\mathbb Q),
\]
where each $\gamma_{0j}$ has nonnegative degree on compact curves.
Fix integers $m_j>0$ such that $m_j\gamma_{0j}$ is the image of an
integral cohomology class. Consider an existing sequence of projective
birational steps
\[
 X_i\xrightarrow{f_i}Z_i\xleftarrow{f_i^+}X_{i+1},
\]
where $f_i^+$ is the identity for a divisorial step, $f_i$ satisfies
Lemma~\ref{prog:integral-h2}, and $\alpha_i$ is trivial on its
contracted ray. Assume every $f_i$-contracted curve has degree
proportional to that ray for the classes in the display. Then the
classes $\gamma_{ij}$ descend through $f_i$ and pull back through
$f_i^+$, the same integers $m_j$ remain integral multiples at every
stage, and the transported classes remain nonnegative on curves.
In particular, for every integral compact curve $C\subset X_i$,
\[
 \alpha_i\cdot C>0\quad\Longrightarrow\quad
 \alpha_i\cdot C\geq\delta,
 \qquad \delta:=\min_j\frac{a_j}{m_j}>0.
\]
\end{lemma}

\begin{proof}
Positivity of the coefficients and nonnegativity of the summands imply
that every $\gamma_{ij}$ annihilates the contracted ray. Apply the
integral descent lemma to the chosen integral lift of
$m_j\gamma_{ij}$ and pull its descended class back to $X_{i+1}$.
This keeps $m_j$ fixed.

To verify nonnegativity, take a curve on $X_{i+1}$. If it is contracted
by $f_i^+$, every descended summand has degree zero. Otherwise its
image is a curve $\Gamma\subset Z_i$. The projective map $f_i$ has a
compact curve mapping onto $\Gamma$ with positive degree: take a
component of its inverse image dominating $\Gamma$ and intersect with
sufficiently many relatively ample hyperplanes. Nonnegativity of the
old summand on this curve implies nonnegativity on $\Gamma$, hence on
the curve on $X_{i+1}$. This argument uses a positive-degree lift, not
a degree-one section. Finally $m_j\gamma_{ij}\cdot C$ is a
nonnegative integer, giving the stated lower bound.
\end{proof}

\begin{corollary}[Uniform crepancy parameter; no termination assertion]
\label{prog:uniform-canonical-parameter}
Suppose additionally that the working spaces are $d$-dimensional klt
spaces, that $\alpha_i$ are nef, and that the ordinary canonical
negative-ray length bound $0<-K_{X_i}\cdot C\leq2d$ is available.
For a fixed real number $t>2d/\delta$, every negative extremal ray
chosen in an existing $(K_{X_i}+t\alpha_i)$-program is
$\alpha_i$-trivial. If the steps are projective ordinary canonical
steps, Lemma~\ref{prog:fixed-integral-grid} therefore applies
inductively with this same value of $t$.

If $\alpha_0$ is not big and $\alpha_0-c_1(K_{X_0})$ is big, then
$K_{X_i}+t\alpha_i$ remains non-pseudo-effective along these crepant
steps. Consequently a terminating such program ends with a Mori
fibre contraction. This corollary does not assert that the requisite
contractions exist or that the program terminates.
\end{corollary}

\begin{proof}
A negative ray is $K_{X_i}$-negative because $\alpha_i$ is nef. Choose
its rational curve generator with the ordinary length bound. If its
$\alpha_i$-degree were positive, then
\[
 (K_{X_i}+t\alpha_i)\cdot C\geq-2d+t\delta>0,
\]
a contradiction. The step is thus an ordinary canonical step with
$\alpha_i$ pulled back from its base, and the preceding lemma applies.
On a common resolution the pullback of
$\alpha_{i+1}-c_1(K_{X_{i+1}})$ is the pullback of
$\alpha_i-c_1(K_{X_i})$ plus the effective canonical comparison.
Thus it stays big. Equality of the crepant pullbacks of $\alpha_i$
preserves its nonbigness. The
identity
\[
 (t+1)\alpha_i=(c_1(K_{X_i})+t\alpha_i)
                  +(\alpha_i-c_1(K_{X_i}))
\]
excludes pseudo-effectivity of its first summand. A nef endpoint is
therefore impossible.
\end{proof}


\begin{lemma}[Cohomology splitting along detected steps]
\label{prog:detected-cohomology}
Let $X_0$ be a smooth compact K\"ahler manifold. Consider a finite
sequence of divisorial contractions and small flips
\[
                         X_0\dashrightarrow X_1\dashrightarrow\cdots
                         \dashrightarrow X_r
\]
between normal globally strongly $\mathbb Q$-factorial compact
K\"ahler spaces with rational singularities.
Assume that every negative contraction $f_i:X_i\to Z_i$ is a
projective contraction of one ray of the analytic cone, that $Z_i$
has rational singularities, and that each small flip is the detected
flip of a global rational line $L_i$ with $L_i\cdot R_i<0$.
For a divisorial contraction assume the usual single exceptional
prime. Then, for every projective smooth resolution $p:U\to X_i$,
\begin{align*}
 H^2(U,\mathbb R)
   &=p^*H^2(X_i,\mathbb R)
       \oplus\bigoplus_{E\text{ exceptional for }p}\mathbb R[E],\\
 H^{1,1}(U,\mathbb R)
   &=p^*H^{1,1}_{\rm BC}(X_i,\mathbb R)
       \oplus\bigoplus_{E\text{ exceptional for }p}\mathbb R[E].
\end{align*}
The first equality also holds over $\mathbb Q$. There are compatible
forward linear transforms on $H^2(-,\mathbb R)$ and on Bott--Chern
cohomology. They are surjective, are isomorphisms for a small flip,
and preserve rational degree-two classes. On a common resolution the
pullback difference is an actual exceptional divisor class; for a
flip it is exceptional over both spaces.

The exceptional quotient also takes the Chern class of a holomorphic
line bundle on $U$ to a rational holomorphic line class on $X_i$.
\end{lemma}

\begin{proof}
For a modification of a smooth compact K\"ahler manifold the two
splittings are the degree-two modification formula. Suppose that they
hold for $X=X_i$. We use only the following birational descent fact:
for a proper bimeromorphic map between normal compact spaces in
Fujiki's class with rational singularities, pullback is injective on
$H^2(-,\mathbb R)$ and on Bott--Chern cohomology, and its image in
either group consists of the classes having degree zero on every
contracted curve. This is the bimeromorphic case of
\cite[Lemma~2.6(1)]{DH2024}.

First note that a degree-two class $\xi\in H^2(X,\mathbb R)$ has
the same degrees on curves as a Bott--Chern class. Indeed, pull it to
a smooth projective resolution $a:W\to X$ and take the $(1,1)$-part
of its Hodge decomposition. The induction hypothesis writes that
part as $a^*\beta+[F]$, with $F$ exceptional over $X$. On every
$a$-contracted curve, $a^*\xi$ and its $(2,0)$ and $(0,2)$ parts
have degree zero. Hence $F$ has degree zero on every such curve.
Exceptional negativity applied to both signs gives $F=0$.
Projective multisections lifting a curve of $X$ then show that
$\xi$ and $\beta$ have the same degrees on every curve of $X$.

Suppose first that $X\dashrightarrow X^+$ is a small detected flip.
Write $f:X\to Z$ and $f^+:X^+\to Z$, and fix an integral curve
$C$ on its negative ray. For $\xi\in H^2(X,\mathbb R)$ set
\[
                    s=\frac{\xi\cdot C}{c_1(L)\cdot C}.
\]
By the preceding paragraph, $\xi-s c_1(L)$ vanishes on every
$f$-contracted curve. Thus it is $f^*\eta$ for a unique
$\eta\in H^2(Z,\mathbb R)$. Define
\[
                    T\xi=(f^+)^*\eta+s c_1(L^+).
\]
For a Bott--Chern class, the same descent takes place in
Bott--Chern cohomology, so the two transforms agree. The detected
flip construction gives an effective rational divisor $E_L$ on a
common projective resolution $a:W\to X$, $b:W\to X^+$ such that
\[
             a^*c_1(L)-b^*c_1(L^+)=[E_L].
\]
Consequently
\[
                         a^*\xi=b^*T\xi+s[E_L].
\]
If $\xi$ is rational, then $s$ is rational. Since the image of
$f^*$ on real degree-two cohomology is the realification of a
rational subspace, its unique preimage $\eta$ is rational as well.
This proves rationality of $T\xi$.

The sets of prime divisors exceptional for $a$ and $b$ are identical:
the two working spaces are small bimeromorphic. The old splitting
and the displayed pullback identity therefore span
\[
              H^2(W,\mathbb R)
                =b^*H^2(X^+,\mathbb R)
                   +\operatorname{span}_{\mathbb R}\{[E]:E\text{ is }b
                   \text{-exceptional}\},
\]
and similarly in bidegree $(1,1)$. The sum is direct. If a pullback
from $X^+$ equals an exceptional divisor class, that divisor has
degree zero on all $b$-contracted curves, so both signs of negativity
make it zero; pullback injectivity then makes the original class
zero. In particular the dimensions on the two sides agree, $T$ is
injective by the old direct sum, and $T$ is surjective. The rational
version follows either from the same argument or from realification.

For a divisorial contraction $f:X\to X^+$, let $D$ be its
exceptional prime. Its degree on the contracted ray is negative:
otherwise $D$ would be $f$-nef, contrary to exceptional negativity.
For any $\xi$, subtract
\[
                    \frac{\xi\cdot C}{D\cdot C}[D]
\]
and descend the result by the same birational cohomology lemma.
This defines $T\xi$. On a common resolution its pullback difference
is a multiple of the total transform of $D$. The target's exceptional
prime set is the old one together with the strict transform of $D$.
The same spanning and directness argument gives the target splitting
and surjectivity; its kernel is $\mathbb R[D]$. Rationality is again
immediate from the quotient of rational intersection numbers.

These arguments initially prove the splitting on a chosen common
resolution. To obtain it on any projective smooth resolution
$U\to X^+$, dominate that resolution and the chosen one by a smooth
common projective resolution. Apply the smooth degree-two
modification formula upstairs and push down to $U$. An exceptional
divisor pushes either to zero or to a divisor exceptional over $X^+$.
This proves spanning on $U$; negativity gives directness there.

Finally let $\mathcal H$ be a holomorphic line on $U$ and put
$\mathcal Q=(p_*\mathcal H)^{**}$. Global strong
$\mathbb Q$-factoriality gives an $m>0$ for which
$\mathcal M=\mathcal Q^{[m]}$ is invertible. The coherent evaluation
comparison gives
\[
                  \mathcal H^{\otimes m}
                  \simeq p^*\mathcal M\otimes\mathcal O_U(F)
\]
for an integral $p$-exceptional divisor $F$. This comparison is
obtained from the coherent direct image and evaluation map; it does
not assume that $\mathcal H$ has a global meromorphic section.
Dividing first Chern classes by $m$ proves the line assertion.
\end{proof}


\begin{corollary}[Stability under projective extraction]
\label{prog:extraction-cohomology}
Let $r:Y\to X$ be a projective bimeromorphic morphism of normal compact
K\"ahler spaces with rational singularities. Suppose that $Y$ is globally
strongly $\mathbb Q$-factorial and that $X$ has the splitting property in
the preceding lemma. Then $Y$ has that property as well. More precisely,
if $E_1,\ldots,E_s$ are the prime divisors exceptional for $r$, then
\[
 H^2(Y,\mathbb R)
   =r^*H^2(X,\mathbb R)\oplus
      \bigoplus_{j=1}^{s}\mathbb R[E_j],
\]
and the analogous statements hold over $\mathbb Q$ and in
Bott--Chern cohomology.
\end{corollary}

\begin{proof}
Choose any projective smooth resolution $q:W\to Y$. The composite
$p=rq$ is a projective resolution of $X$. Its exceptional primes are
exactly the $q$-exceptional primes together with the strict transforms
$\widetilde E_j$ of the $E_j$. Since every $E_j$ is $\mathbb Q$-Cartier,
\[
 [\widetilde E_j]=q^*[E_j]+[F_j]
\]
for a rational $q$-exceptional divisor $F_j$. Substitute these identities
in the splitting for $p$. The result spans $H^2(W,\mathbb R)$ by
$q^*H^2(Y,\mathbb R)$ and the $q$-exceptional divisor classes. Their sum
is direct: an exceptional divisor whose class is a pullback has degree
zero on every $q$-contracted curve, and applying exceptional negativity
to both signs makes the divisor zero. Pullback injectivity then applies.
The same argument in bidegree $(1,1)$ and over $\mathbb Q$ proves the
splitting property for $Y$.

Now expand $q^*\xi$, for $\xi\in H^2(Y,\mathbb R)$, using the original
$p$-splitting and replace the $[\widetilde E_j]$ as above. The difference
between $q^*\xi$ and a class in
$q^*(r^*H^2(X,\mathbb R)+\sum_j\mathbb R[E_j])$ is $q$-exceptional,
and is therefore zero by the direct sum just proved. Thus the displayed
formula for $H^2(Y,\mathbb R)$ spans. Its directness follows from
exceptional negativity for $r$ and pullback injectivity. The rational
and Bott--Chern statements follow by the identical argument. The line
bundle assertion of the preceding lemma uses only global strong
$\mathbb Q$-factoriality and hence applies to $Y$ as well.
\end{proof}



\paragraph{Support in the special-termination comparison.}
\label{prog:special-support}
The following support comparison applies to the detected steps in
the dimension induction. Let $\psi:X\dashrightarrow X^+$ be a detected
small negative step with maps $f:X\to Z$ and $f^+:X^+\to Z$,
and let $L$ be its negative detector. For a sufficiently
divisible $r>0$, put
\[
 \operatorname{im}(f^*f_*L^r\longrightarrow L^r)
      =\mathcal I\otimes L^r.
\]
The zero set of $\mathcal I$ is exactly the exceptional locus of $f$.
Outside that locus $f$ is an isomorphism, so the evaluation map is
surjective. On a positive-dimensional projective fibre, every section
of $L^r$ vanishes: its restriction to each integral curve has negative
degree, and those curves cover every positive-dimensional fibre
component. Connectedness supplies the whole nontrivial fibre.

On a resolution principalizing $\mathcal I$, the detector comparison
is a positive multiple of the effective divisor defined by
$\mathcal I\mathcal O_W$. Consequently its support is the full
inverse image of the exceptional locus. The driving adjoint differs
from a positive real multiple of $c_1(L)$ by a class pulled back from
the contraction base. Its comparison divisor is therefore the same
positive multiple of the detector comparison: their class difference
is exceptional and zero, so both-sign negativity makes it zero as a
divisor. The comparison for a longer
negative program is at least this first-step divisor, by discrepancy
monotonicity on a common resolution.

Suppose the beginning and end restrictions to a surviving normalized
lc stratum $S$ are isomorphic and their strict adjunction data agree.
The general point of this stratum avoids every intervening surgery:
a zero-discrepancy place remains such a place, whereas strict
comparison would increase its discrepancy if its center were contained
in the exceptional locus. Thus its strict transform $S_W$ on a common
resolution is not contained in the comparison support. Iterated strict
adjunction identifies the restriction of the effective comparison with
the difference of the identified adjoints, so its class is zero.
A nonzero effective divisor on a positive-dimensional compact K\"ahler
stratum has positive K\"ahler mass. The restricted divisor is therefore
zero. If the first step meeting the stratum were nontrivial there,
the full inverse-image support and surjectivity $S_W\to S$ would make
this restriction nonzero, a contradiction. For a zero-dimensional
stratum, generic-point avoidance already excludes an intersection.
The whole program is consequently an isomorphism near the stratum.
This is the same strict-adjunction support argument used in the proof
of Theorem~\ref{prog:special-termination}; it does not use its
lower-dimensional abundance assumption.

\begin{proposition}[A uniform relative NQC bound]
\label{prog:relative-nqc-grid}
Let $f_0:U\to Z$ be an already projective surjective morphism with
connected fibres between smooth compact
K\"ahler manifolds, and suppose that
\[
 f_0^*:H^0(Z,\Omega_Z^2)\longrightarrow H^0(U,\Omega_U^2)
\]
is an isomorphism. Let $\alpha_0$ be a nef class with an expression
\[
       \alpha_0=\sum_{j=1}^k r_j\eta_{j,0},\qquad
       r_j>0,\quad \eta_{j,0}\in H^2(U,\mathbb Q),
\]
where $\eta_{j,0}$ has nonnegative degree on every curve vertical
over $Z$. Let $D_0=K_U+\Delta_U+\mathbf M_U$ be a generalized klt
adjoint. Consider its projective relative program for
$D_0+a\alpha_0$, under the projective local ordinary reduction described in the
proof of Lemma~\ref{prog:base-point-free-induction}.

If $\alpha_0=0$, triviality is immediate. Otherwise there are
positive integers $m_j$, chosen once on $U$, and a number
\[
                         \delta=\min_j\frac{r_j}{m_j}>0
\]
such that, for $a\delta>2\dim U$, the entire relative program is
$\alpha_0$-trivial. The same number $a$ works at every step.
\end{proposition}

\begin{proof}
\emph{Rational classes modulo the fixed base.}
Choose an $f_0$-ample line bundle with integral Chern class $H$, and
put $e=\dim U-\dim Z$ and $c=f_{0*}(H^e)>0$. The rational operator
\[
 \Pi(\eta)=\eta-
 f_0^*\!\left(\frac{f_{0*}(\eta\smile H^e)}{c}\right)
\]
preserves type $(1,1)$ and kills types $(2,0)$ and $(0,2)$.
Indeed these latter classes are pulled back from $Z$, and the
projection formula applies. Thus
$\Pi(H^2(U,\mathbb Q))\subset H^{1,1}(U)\cap H^2(U,\mathbb Q)$.
By the Lefschetz $(1,1)$ theorem, choose an actual line bundle
$L_{j,0}$ and an integer $m_j>0$ with
\[
 \eta_{j,0}=\frac1{m_j}c_1(L_{j,0})+f_0^*\gamma_j,
 \qquad \gamma_j\in H^2(Z,\mathbb Q).
\]
The original summands $\eta_{j,0}$ need not have type $(1,1)$.
Their base components account for this. Their weighted sum gives
\begin{equation}\label{prog:eq:relative-nqc}
 \alpha_0=\sum_j\frac{r_j}{m_j}c_1(L_{j,0})+f_0^*\Gamma,
 \qquad \Gamma=\sum_jr_j\gamma_j\in H^{1,1}(Z,\mathbb R).
\end{equation}
The final assertion follows because $f_0^*$ is an injective Hodge
map and the other terms have type $(1,1)$. The identity is therefore
also one of Bott--Chern classes.
Each $L_{j,0}$ has nonnegative integral degree on every
$Z$-vertical curve.

\smallskip\noindent
\emph{Induction across an actual contraction.}
Suppose at a finite stage $f_i:U_i\to Z$ we have actual line bundles
$L_{j,i}$, nef over $Z$, and
\[
 \alpha_i=\sum_j\frac{r_j}{m_j}c_1(L_{j,i})+f_i^*\Gamma.
\]
Assume also that $\alpha_i$ is nef and that the transformed
$D_i$ is generalized klt. These assertions hold initially.
Let $R$ be a $(D_i+a\alpha_i)$-negative extremal ray over $Z$.
It is $D_i$-negative because $\alpha_i$ is nef. Lemma~\ref{prog:relative-hodge} and Lemma~\ref{prog:projective-cone}
identify the relative curve cone with the corresponding face of the
full analytic cone, at the initial model and after every step. The
analytic cone theorem therefore supplies a rational generator $C$ with
\[
                     0<-D_i\cdot C\leq2\dim U.
\]
All $L_{j,i}\cdot C$ are nonnegative integers. If
$\alpha_i\cdot C>0$, one of them is at least one, and hence
$\alpha_i\cdot C\geq\delta$. Therefore
\[
 (D_i+a\alpha_i)\cdot C
             \geq-2\dim U+a\delta>0,
\]
a contradiction. We have $\alpha_i\cdot C=0$ and
$L_{j,i}\cdot C=0$ for every $j$. Every curve contracted by the
projective extremal contraction $h_i:U_i\to Y_i$ has class on
$R$ when tested by global line bundles, so the same vanishing
holds for all of them.

On each fixed Stein chart of $Z$, choose the effective real
ordinary representative of the initial relatively ample nef
datum once, before running the program. Its actual real linear
equivalence to the fixed real line-bundle representative persists
under pushforward. The Bott--Chern comparison with the generalized
adjoint modulo the fixed base class persists separately by
exceptional negativity, as detailed below.
Thus every $h_i$ is, locally on $Y_i$, a contraction with an
ordinary real klt log-Fano boundary. On a relatively compact base
chart, rational approximation in the finite affine system of Cartier
adjoint identities gives a rational effective klt boundary with the
same antiample sign, as in Proposition~\ref{prog:detected-flip}.
Lemma~\ref{prog:integral-descent} therefore applies.
Consequently
\[
 M_{j,i}:=(h_i)_*L_{j,i}\ \text{is a line bundle},\qquad
 h_i^*M_{j,i}=L_{j,i}.
\]
For a divisorial contraction set $L_{j,i+1}=M_{j,i}$.
For a flip $h_i^+:U_{i+1}\to Y_i$, set
$L_{j,i+1}=(h_i^+)^*M_{j,i}$.
These are actual line bundles with the \emph{same} integers $m_j$.
The descended Bott--Chern class
\[
 \alpha_{Y_i}
     =\sum_j\frac{r_j}{m_j}c_1(M_{j,i})+f_{Y_i}^*\Gamma
\]
satisfies $h_i^*\alpha_{Y_i}=\alpha_i$.
It is nef by descent of nefness through the projective surjection
$h_i$; its pullback to the next model is $\alpha_{i+1}$.
This proves that the step is crepant for $\alpha_i$. Hence it is
also a $D_i$-negative step, and the original generalized pair
remains gklt.

Finally each $M_{j,i}$ is nef over $Z$. For a $Z$-vertical curve
$B\subset Y_i$, projectivity supplies a curve $B'\subset U_i$
mapping onto it with some positive degree $d_B$. Then
\[
 d_B(M_{j,i}\cdot B)=L_{j,i}\cdot B'\geq0.
\]
No assertion $d_B=1$ is required. Pullback preserves this relative
nefness, so the induction applies on $U_{i+1}$. Its curve degrees
are again integers because the transported objects are line
bundles, not because any curves were lifted with degree one.
This proves the fixed bound through the whole program.
\end{proof}

\begin{remark}[Scope of the degree grid]
The transported rational classes
\[
 \eta_{j,i}=\frac1{m_j}c_1(L_{j,i})+f_i^*\gamma_j
\]
have degrees in $m_j^{-1}\mathbb Z_{\geq0}$ on curves vertical over
the fixed smooth base $Z$. This relative assertion is exactly what
the initial projective program requires. It does not assert
nonnegativity of these individual classes on all curves of every
birational model, and uses neither rational connectedness of
resolution fibres nor Graber--Harris--Starr.
\end{remark}

\subsubsection{The negative-part comparison for a nef adjoint}
We write $N_\sigma(\xi)$ for the divisorial negative part of a
pseudo-effective class. We use its homogeneity, subadditivity,
continuity after adding a vanishing K\"ahler perturbation, and the
identities
\[
 N_\sigma(r^*\xi+[E])=N_\sigma(r^*\xi)+E,
 \qquad r_*N_\sigma(r^*\xi)=N_\sigma(\xi)
\]
for an effective $r$-exceptional divisor $E$. A nef class has zero
negative part. These are the divisorial Zariski-decomposition
properties of \cite[Section~3]{Boucksom2004} and
\cite[Appendix~A]{DHY}. No identity
$N_\sigma(r^*\xi)=r^*N_\sigma(\xi)$ for arbitrary $\xi$ is assumed.

\begin{lemma}[The absolute negative part after a relative program]
\label{prog:negative-part-program}
Let $\mu:U\to X$ be a projective resolution, let $\alpha$ be nef on $X$,
let $c>0$, and let $F_U\geq0$ be $\mu$-exceptional. Put
\[
 D_U=c\mu^*\alpha+[F_U].
\]
Suppose that $\phi:U\dasharrow V$ is a finite $D_U$-negative program,
possibly relative to another base, and denote its transformed class and
divisor by $D_V$ and $F_V=\phi_*F_U$. Then
\[
                         N_\sigma(D_V)=F_V.
\]
\end{lemma}

\begin{proof}
Take a common projective resolution $p:W\to U$, $q:W\to V$.
The negativity comparison for the finite program is
\[
                p^*D_U=q^*D_V+[E],
 \qquad E\geq0,\quad E\text{ is }q\text{-exceptional}.
\]
This comparison is valid for a relative negative program: its proof
uses the negativity of the contracted rays and the positive adjoint on
each flipped side, not absolute nefness of the final adjoint.
Since $p^*F_U$ is effective and exceptional over $X$, the exceptional
identity and nefness of $(\mu p)^*\alpha$ give
\[
 N_\sigma(p^*D_U)=p^*F_U
                 =N_\sigma(q^*D_V)+E.
\]
Pushing forward by $q$ proves the assertion. In particular, the
negative part in this statement is absolute, even when the program
that produced $V$ was relative.
\end{proof}

\subsubsection{Descent of the effective vertical divisor}

\begin{lemma}[Vertical numerical triviality]
\label{prog:vertical-divisor}
Let $g:V\to S$ be a projective surjective morphism with connected fibres.
Assume that $S$ is globally Weil $\mathbb{Q}$-factorial. Let $F\geq0$ be
an effective real Cartier divisor on $V$, vertical over $S$, such that
\[
                    F\cdot C=0
       \quad\text{for every curve }C\text{ contracted by }g.
\]
Then there is an effective real Cartier divisor $F_S$ on $S$ such that,
as actual divisors,
\[
                              F=g^*F_S.
\]
\end{lemma}

\begin{proof}
For a prime divisor $P\subset S$, write $m_Q$ for the multiplicity of
$Q$ in $g^*P$, where $Q$ ranges over the prime divisors dominating $P$.
These multiplicities may be computed over the smooth generic locus of
$P$, where $P$ is Cartier. Put
\[
 b_P=\min_{g(Q)=P}\frac{\operatorname{coeff}_Q F}{m_Q},
 \qquad F_S=\sum_P b_PP.
\]
Only finitely many $b_P$ are nonzero, since any such $P$ is the image of
a component of $F$. Thus $F_S$ is a genuine effective Weil real divisor.
Global Weil $\mathbb{Q}$-factoriality makes this finite divisor real
Cartier.

We first check equality over codimension one in $S$. Work near a
general point of $P$ and restrict to a transverse disk. Resolving the
source, and cutting by general relative ample hypersurfaces if the
relative dimension exceeds one, reduces to a projective surface over
that disk with connected fibres. These operations can be performed
over a relatively compact Stein neighbourhood; they do not require
global sections on $S$. The intersection matrix of the components of
the special fibre is negative semidefinite, with kernel generated by
the full fibre with its multiplicities. The restriction of $F$ has
zero intersection with every fibre component. Its coefficients are
therefore proportional to those multiplicities. Equivalently, all the
ratios in the definition of $b_P$ are equal. When $g$ is birational this
assertion over the generic point of $P$ is immediate.

Consequently
\[
                              G=F-g^*F_S
\]
is a signed real Cartier divisor supported over a subset of codimension
at least two in $S$. Moreover, $G$ is numerically trivial over $S$.
We check that such a signed exceptional divisor is zero. The assertion
is local on $S$. Over a relatively compact Stein neighbourhood, take
$d=\dim V-\dim S$ general relative ample hypersurfaces and resolve their
intersection. They give a projective generically finite morphism
$H\to S$. The cuts may be chosen to meet any prescribed component of
$G$ in a nonzero divisorial trace. After normalization and Stein
factorization, $H\to S$ factors as a projective birational morphism
$H\to S'$ followed by a finite morphism $S'\to S$. The restricted
divisor $G|_H$ is exceptional for $H\to S'$ and numerically trivial
over $S'$. Applying the ordinary exceptional negativity lemma to both
$G|_H$ and $-G|_H$ gives $G|_H=0$. The choice of the cuts therefore
excludes every nonzero component of $G$. Hence $G=0$.

The surface argument above is also the usual proof of the
degenerate-divisor negativity statement: after subtracting the minimum
multiple of the full fibre, an effective residual fibre divisor omits
a component and cannot be numerically trivial. Notice that projectivity
of $g$ was a hypothesis; it was not deduced from factoriality.
\end{proof}

\subsubsection{The negative part on the lower-dimensional base}

\begin{lemma}[Detecting the negative part by pullback currents]
\label{prog:negative-part-base}
Let $g:V\to S$ be a surjective morphism. Suppose that $\alpha_S$ is nef,
$F_S\geq0$ is real Cartier, $c>0$, and
\[
 D_S=c\alpha_S+[F_S],\qquad D_V=g^*D_S,\qquad
 N_\sigma(D_V)=g^*F_S.
\]
Then $N_\sigma(D_S)=F_S$.
\end{lemma}

\begin{proof}
Nefness gives $N_\sigma(D_S)\leq F_S$. Fix a prime divisor $P\subset S$
and a prime divisor $Q\subset V$ dominating it. Write
$m=\operatorname{mult}_Q(g^*P)>0$ and $b=\operatorname{coeff}_P F_S$.
At their generic smooth points, pullback of positive currents satisfies
\[
                         \nu(g^*T,Q)=m\nu(T,P).
\]
Indeed, the divisorial term $\nu(T,P)[P]$ pulls back with multiplicity
$m$, and the residual current has zero generic Lelong number there.

Choose K\"ahler forms $\omega_S,\omega_V$ and a constant $C>0$ for which
$C\omega_V-g^*\omega_S$ is K\"ahler. For $\varepsilon>0$, let $T_\varepsilon$
be any positive current in the big class $D_S+\varepsilon[\omega_S]$.
Minimality of the multiplicity and monotonicity under adding a nef
class give
\begin{align*}
 m\nu(T_\varepsilon,P)
 &=\nu(g^*T_\varepsilon,Q)\\
 &\geq\nu(D_V+\varepsilon g^*[\omega_S],Q)\\
 &\geq\nu(D_V+C\varepsilon[\omega_V],Q).
\end{align*}
Taking the infimum over $T_\varepsilon$, and then letting
$\varepsilon\downarrow0$, gives
\[
 m\nu(D_S,P)\geq\nu(D_V,Q)
                 =\operatorname{coeff}_Q(g^*F_S)=mb.
\]
This proves the reverse inequality for every $P$.
All multiplicities are unchanged by resolving away from the generic
points in question, so the same argument applies to the normal spaces
under consideration.

The $\varepsilon$ perturbation is necessary: at a pseudoeffective
boundary class, the infimum over its exact positive currents need not
equal its minimal multiplicity. The argument uses that equality only
for the big perturbed classes.
\end{proof}

\begin{lemma}[Pullback when the positive part is nef]
\label{prog:negative-part-pullback}
Suppose
\[
 \xi=P+[F],\qquad P\text{ nef},\quad F\geq0\text{ real Cartier},
 \qquad N_\sigma(\xi)=F.
\]
For every projective resolution $p:W\to S$,
\[
                         N_\sigma(p^*\xi)=p^*F.
\]
\end{lemma}

\begin{proof}
Write $N'=N_\sigma(p^*\xi)$. Since $p^*P$ is nef, $N'\leq p^*F$.
Birational invariance gives $p_*N'=F$, so
$E=p^*F-N'\geq0$ is $p$-exceptional. The positive part of $p^*\xi$ is
\[
                         p^*P+[E],
\]
and is modified nef. If $E\neq0$, exceptional negativity in its
covering-curve form provides a component of $E$ covered by curves $C_t$
contracted by $p$, with $E\cdot C_t<0$
\cite[Appendix~A, Lemma~A.3]{DHY}. A modified nef class has nonnegative
intersection with a general member of a family of curves covering a
prime divisor: use currents with arbitrarily small negative part and
zero generic divisorial Lelong number, restrict to a general member,
and let the negative bound tend to zero. But here
\[
                       (p^*P+E)\cdot C_t=E\cdot C_t<0,
\]
a contradiction. Thus $E=0$.
\end{proof}

\begin{proposition}[Any lower-dimensional good model suffices]
\label{prog:nef-good-comparison}
Suppose the initial relative program has the data in
Lemma~\ref{prog:negative-part-program}. Assume in addition that
its class $\alpha_V$ is nef and agrees with $\mu^*\alpha$ on a common
resolution, and that there is an already projective connected-fibre
morphism $g:V\to S$ with
\[
\alpha_V=g^*\alpha_S,\qquad D_V=g^*D_S.
\]
Suppose there is a projective small modification
$s:S^{\mathrm q}\to S$ such that $S^{\mathrm q}$ is globally Weil
$\mathbb Q$-factorial; the identity is allowed. Assume that the
lower-dimensional adjoint $s^*D_S$ has a good model:
there are a normal compact K\"ahler space $S_m$, a common projective
resolution $p:W\to S^{\mathrm q}$, $q:W\to S_m$, and a nef semiample
class $D_m$ such that
\[
                  p^*s^*D_S=q^*D_m+[E],
 \qquad E\geq0,\quad E\text{ is }q\text{-exceptional}.
\]
Then $\alpha$ is semiample on $X$. If the contraction defining
semiampleness of $D_m$ is Moishezon, the resulting contraction of
$\alpha$ is Moishezon as well. In particular, it is unnecessary to
lift a program on $S$ to $V$, or to require that a chosen program
producing $S_m$ preserve $\alpha_S$ at every intermediate step.
\end{proposition}

\begin{proof}
Nefness descends under the surjective morphism $g$, so $\alpha_S$ is nef.
The class of $F_V$ is pulled back from $S$, so $F_V$ is numerically
trivial over $S$. It is therefore vertical: an effective horizontal
component would restrict to a nonzero effective divisor on a general
projective fibre, with positive intersection against a suitable power
of an ample class, contradicting numerical triviality. For a
birational $g$, verticality is automatic by dimension.
First take a projective resolution $\pi:\widetilde V\to V$ of the main
component of $V\times_S S^{\mathrm q}$. The induced map
$\widetilde g:\widetilde V\to S^{\mathrm q}$ is projective and has
connected fibres. Pull back $\alpha_V,D_V,F_V$ along $\pi$, and pull
back $\alpha_S,D_S$ along $s$. All displayed class identities are
preserved, and $\pi^*F_V$ remains an effective vertical divisor.
Lemma~\ref{prog:negative-part-program} gives
$N_\sigma(D_V)=F_V$ before this modification. Because
$D_V=c\alpha_V+[F_V]$ with $\alpha_V$ nef,
Lemma~\ref{prog:negative-part-pullback} then gives
\[
              N_\sigma(\pi^*D_V)=\pi^*F_V.
\]
We may therefore replace $(V,S,g)$ by
$(\widetilde V,S^{\mathrm q},\widetilde g)$ and suppress the new
superscripts in the rest of the proof. The good-model comparison now
reads $p^*D_S=q^*D_m+[E]$. No assertion that the crepant subboundary on
$\widetilde V$ is effective is needed: this space is used only for
classes, currents and the effective divisor $\pi^*F_V$. The
lower-dimensional generalized pair is the crepant pullback to the
small model $S^{\mathrm q}$.

The class identity $[F_V]=g^*(D_S-c\alpha_S)$ shows that $F_V$ is
numerically trivial over $S$. Lemma~\ref{prog:vertical-divisor}
gives an actual effective real Cartier divisor $F_S$ with
$F_V=g^*F_S$. Injectivity of pullback yields
\[
                         D_S=c\alpha_S+[F_S].
\]
The established identity $N_\sigma(D_V)=F_V$ and
Lemma~\ref{prog:negative-part-base} therefore give
$N_\sigma(D_S)=F_S$, and
Lemma~\ref{prog:negative-part-pullback} gives
\[
                         N_\sigma(p^*D_S)=p^*F_S.
\]
On the other hand $q^*D_m$ is nef, so the exceptional-divisor identity
applied to the good-model comparison gives
\[
                         N_\sigma(p^*D_S)=E.
\]
Thus $E=p^*F_S$ as actual divisors. Subtracting them from the comparison
proves the exact class equality
\[
                         q^*D_m=c\,p^*\alpha_S.
\]
This proves the needed preservation of $\alpha_S$ from the final good
model, without an assumption about its intermediate models.

Choose a contraction $h:S_m\to T$ to a normal compact K\"ahler space
and a K\"ahler class $\kappa$ on $T$ with $D_m=h^*\kappa$. Taking a
resolution of the main component of $V\times_S W$, and then a common
projective resolution with $U$, produces a projective modification
$r:R\to X$ and a holomorphic map $\ell:R\to T$ satisfying
\[
                         r^*\alpha=\frac1c\ell^*\kappa.
\]
For every curve $C$ in a fibre of $r$ this equality implies
$\ell(C)$ is a point, since a K\"ahler class has positive degree on
every nonconstant image curve. The fibres of the projective modification
$r$ are connected projective complex spaces, hence are connected by
chains of curves. Therefore $\ell$ is
constant on each fibre of $r$. The factorization theorem for a proper
map onto a normal space gives a holomorphic map $f:X\to T$ with
$\ell=f\circ r$. Pushing forward the class equality by $r$ gives
\[
                              \alpha=f^*(\kappa/c).
\]
The maps from the main fibre-product component to $S_m$ have connected
general fibres; their Stein factorizations are finite birational over
the normal space $S_m$, hence have connected fibres everywhere.
Together with the connected fibres of $h$, this shows that $\ell$,
and therefore $f$, has connected fibres. This is the required
semiampleness of $\alpha$.

Finally, $R\to S_m$ is projective: it is obtained from the projective
map $g$ by base change, followed by projective resolutions and the
projective map $q$. If $h$ is Moishezon, choose a projective surjection
$H\to S_m$ such that $H\to T$ is projective. A component of
$R\times_{S_m}H$ dominating $R$ is projective and surjective over $X$,
and its map to $T$ is projective. This is a Moishezon witness for $f$.
\end{proof}


\begin{lemma}[The contraction step in the dimension induction]
\label{prog:base-point-free-induction}
Assume $\mathsf C_{d-1}$, $\mathsf B_{d-1}$ and $\mathsf M_{d-1}$.
Then $\mathsf C_{d,\mathrm{big}}$ and $\mathsf B_d$ hold.
\end{lemma}
\begin{proof}
\smallskip\noindent
\emph{The big non-klt contraction.}
Apply the construction of \cite[Section~3]{HX}. On its smooth
modification $\nu:U\to X$ it writes
\[
 \nu^*\alpha=[D_U]+\eta_U,
 \qquad D_U\geq0,\quad \eta_U\text{ K\"ahler}.
\]
The induction over the jumping coefficients of the multiplier
ideal reduces the new reduced stratum to dimension at most $d-1$.
The given contraction on the old non-klt subspace and
$\mathsf C_{d-1}$ provide the Moishezon maps of those strata.
Every map that is promoted to a projective map in that construction
contracts exactly the curves on which $\nu^*\alpha$ vanishes.
On each such curve the actual real line bundle $-D_U$ has degree
$\eta_U\cdot C$. On a further projective common resolution
$q:U''\to U$, choose an effective exceptional divisor $E$ with
$-E$ relatively ample and choose $\varepsilon>0$ sufficiently small.
Use the decomposition
\[
 q^*\nu^*\alpha=[D'']+\eta'',\qquad
 D''=q^*D_U+\varepsilon E,\qquad
 \eta''=q^*\eta_U-\varepsilon[E].
\]
Here $\eta''$ is K\"ahler, and the actual real line bundle $-D''$
has degree $\eta''\cdot C$ on every contracted curve. Restrict this
decomposition to the smooth reduced components in the gluing
construction. Lemma~\ref{prog:actual-line-polarization} then supplies
precisely the relative ampleness needed in place of
\cite[Lemma~2.42]{HX}, including the common-resolution step in its
Claim~3.5. The exceptional term $-\varepsilon[E]$ makes the
positive part $\eta''$ K\"ahler on the common resolution.

The exact sequences of multiplier ideals, relative vanishing,
finite pushouts, and extension from sufficiently thickened
$\operatorname{Supp}D_U$ in that proof then apply unchanged.
They give the birational Moishezon contraction in
$\mathsf C_{d,\mathrm{big}}$. In particular the gluing step uses
neither generalized termination nor a projectivity criterion for
an undetected ray.

\smallskip\noindent
\emph{Preparation for both cases of $\mathsf B_d$.}
The modified-big perturbation
\cite[Lemma~2.23]{HX} supplies a nef datum that dominates a
K\"ahler class on its carrier. Subtracting a sufficiently small
multiple of the pullback of a K\"ahler class $\omega_X$ still leaves
a nef datum on that carrier. Thus
$\alpha$ is a gklt adjoint plus $\varepsilon\omega_X$, and the cone
argument of \cite[Lemma~2.45]{HX} makes $\alpha$ NQC.
This preparation applies whether or not $\alpha$ is big.

\smallskip\noindent
\emph{The nef-and-big case of $\mathsf B_d$.}
For a gklt pair the multiplier ideal is the unit ideal, so the
preceding big non-klt assertion now applies and gives a birational
contraction $h:X\to Y$.
On a projective resolution $p:W\to X$ chosen projective over $Y$,
relative Kawamata--Viehweg vanishing for the effective exceptional
divisor $-\lfloor B_W\rfloor$ gives
$R^ih_*\mathcal O_X=0$ for $i>0$. Thus $Y$ has rational
singularities. Proper birational Bott--Chern descent
\cite[Lemma~8.7]{DHP} gives $\alpha=h^*\gamma$.
The descended adjoint is gklt, nef and big, and has no trivial
curve. The criterion \cite[Theorem~4.3]{HLX} makes $\gamma$
K\"ahler.

For completeness, the projectivity input in that criterion is
restricted to a map $S\to Z'$ between smooth compact K\"ahler
manifolds with rationally connected general fibre: $S$ is the
chosen smooth divisor on a resolution and $Z'$ is the resolution
of the null-locus component. The smooth-source, smooth-base
argument in \cite[Theorem~3.1, Step~1]{CH} applies. Its subsequent
relative program starts with this projective map. Thus this application of the criterion uses only the
smooth-source, smooth-base projectivity argument and an
already-projective relative program.

\smallskip\noindent
\emph{The non-big case of $\mathsf B_d$.}
Use a projective small factorialization and the modified-big
perturbation of the pair. The non-pseudo-effectivity of $K_X$
gives an MRC fibration by \cite{Ou2025}. Choose a smooth K\"ahler
MRC base $Z$ and a smooth modification $\mu:U\to X$ such that
$U\to Z$ is projective. Only the smooth case of
\cite[Theorem~3.1, Step~1]{CH} is needed: pullback identifies
the holomorphic two-forms since the general fibre is rationally
connected. Lemma~\ref{prog:relative-hodge} identifies any $(1,1)$-class on $U$,
modulo a class pulled back from $Z$, with an actual real line bundle.
Lemma~\ref{prog:projective-cone} identifies its relative rays
with rays of the analytic cone; both statements persist along the
projective relative construction.

Write, as in \cite[Theorem~4.1, Claim~4.1]{HX},
\[
 D_U(a)=K_U+\Delta_U+\mathbf M_U+a\alpha_U
     =(a+1)\alpha_U+F_U,
 \quad \alpha_U=\mu^*\alpha,\quad
 \Delta_U,F_U\geq0,
\]
where $\Delta_U$ is klt and $F_U$ is $\mu$-exceptional.
The modified-big replacement is chosen, as in the cited Claim~4.1,
so that the nef datum $\mathbf M_U$ on this smooth carrier is
K\"ahler. If a further projective resolution is needed, subtract a
sufficiently small effective exceptional divisor from its nef part
and add that divisor to the subboundary; this preserves the
adjoint and the klt inequalities.
Choose $a_0$ using the uniform bound of
Proposition~\ref{prog:relative-nqc-grid}. It makes the relative
$D_U(a_0)$-program $\alpha_U$-trivial and preserves that choice
through every step. This is a
\emph{projective} relative program from its first step: the nef
data are represented by real line bundles modulo $Z$.
Lemma~\ref{prog:relative-hodge} applies on every
intermediate model over this fixed smooth base.
More precisely, put $r:U\to Z$ and choose one actual global real
line bundle $L$ with
\[
 \mathbf M_U+a_0\alpha_U=c_1(L)+r^*\gamma.
\]
The bundle $L$ is relatively ample, because $\mathbf M_U$ is
K\"ahler and $\alpha_U$ is nef. On each of a fixed finite collection
of relatively compact Stein charts of $Z$, choose an effective
real divisor $\Theta$ representing $L$, with $(U,\Delta_U+\Theta)$
klt. The representatives can retain a positive relatively ample
part in their boundaries. Here the equality
$K_U+\Delta_U+\Theta\sim_{\R}K_U+\Delta_U+L$ is an actual
real linear equivalence of line bundles; the equality with
$D_U(a_0)$ modulo $r^*\gamma$ is separately an equality of
Bott--Chern classes.

Fix these ordinary pairs on the initial charts. Their actual
linear equivalences and their Bott--Chern comparisons with $D_U(a_0)$
modulo the fixed base class persist separately throughout a global
relative program. Proposition~\ref{prog:projective-good-model} applies:
the source is smooth, its adjoint is represented by an actual global
real line modulo $Z$, its carrier nef part is relatively ample, and
$D_U(a_0)=(a_0+1)\mu^*\alpha+[F_U]$ is pseudo-effective.
It gives a finite program $U\dasharrow V$ and a good relative endpoint.
This is one global program controlled by fixed ordinary chart pairs.
Its construction uses global line algebras, not a choice of unrelated
local minimal models. Proposition~\ref{prog:relative-nqc-grid} makes
every step $\alpha_U$-trivial with the same value $a_0$.

Here is an explicit ample-model comparison that gives the required
descent of $\alpha_V$. Both $D_V(a_0)$ and $\alpha_V$ are nef over
$Z$. Put $a_1=a_0+1$ and $a_2=a_0+2$. Before the program starts,
choose fixed ordinary representatives on the finite chart cover for
all three relatively ample nef parts $\mathbf M_U+a_j\alpha_U$,
$j=0,1,2$, on one simultaneous log resolution. The
$\alpha_U$-crepancy of the constructed program makes its endpoint
$V$ a weak model for each of these three fixed ordinary adjoints.
Their traces $D_V(a_j)=D_V(a_0)+(a_j-a_0)\alpha_V$ are nef.
Corollary~\ref{prog:projective-semiample} therefore makes
\[
 D_V(a_j)=D_V(a_0)+(a_j-a_0)\alpha_V
\]
semiample over $Z$. Their zero curves are exactly the common zero
curves of $D_V(a_0)$ and $\alpha_V$. Their projective ample-model
contractions therefore have the same fibres: these fibres are
connected by curves, and each contraction is constant on the
fibres of the other. Identify the two normal targets, writing the
common contraction as $g:V\to S$. If
$D_V(a_j)=g^*\lambda_j$, with $\lambda_j$ relatively K\"ahler over
$Z$, subtraction gives
\[
 \alpha_V=g^*(\lambda_2-\lambda_1)=g^*\alpha_S.
\]
This is an equality of Bott--Chern classes, not merely of curve
degrees. Since the initial program is $\alpha_U$-trivial, it is
also $D_U(a_1)$-negative. From now on put $a=a_1$ and suppress
the argument in $D_U(a)$ and $D_V(a)$.

The maps $V\to Z$ and $S\to Z$ are projective. Moreover $g$ is
projective: a relatively ample bundle for $V\to Z$ restricts to
an ample bundle on every fibre of $g$, and hence is $g$-ample.
Each $D_U(a_j)$ is not big globally: pushing a big such class to
$X$ would make $(a_j+1)\alpha$ big. If it were big over the smooth
MRC base $Z$, relative-canonical positivity
\cite[Theorem~2.48]{HX}, together with the pseudo-effectivity of
$K_Z$, would make it big globally. It is therefore not big over
$Z$. This property persists to its relative semiample model, so
its relative canonical morphism has $\dim S<d$.
Apply Proposition~\ref{prog:base-presentation} to $g$ and the
whole adjoint $D_V$. Its generalized data are effective and gklt.
The nef datum is represented on the initial carrier by
$\mathbf M_U+a\alpha_U$, a K\"ahler class; on a common carrier over
$V$ its pullback is nef and big. Adding the strict transform of the
effective boundary proves that the total boundary on $V$ is globally
modified big. The proposition therefore gives
\[
 D_S=K_S+B_S+\mathbf N_S,\qquad D_V=g^*D_S,
\]
with effective gklt data on $S$, a globally nef carrier, and globally
modified-big total boundary. Independently, $\alpha_S$ is nef by
descent of nefness through the surjective projective map $g$.
The lower-dimensional model assertion will be applied to this whole
adjoint $D_S$.

Take a projective small factorialization $s:S^{\mathrm q}\to S$
and resolve the main component of $V\times_S S^{\mathrm q}$.
The resulting maps $\pi:\widetilde V\to V$ and
$\widetilde g:\widetilde V\to S^{\mathrm q}$ are projective,
and $\widetilde V$ is compact K\"ahler. Pull back $D_V$,
$\alpha_V$ and $F_V$ to $\widetilde V$, and $D_S$ and
$\alpha_S$ to $S^{\mathrm q}$. The pair on $S^{\mathrm q}$ is
gklt and its boundary-plus-nef class remains modified big.
Lemma~\ref{prog:negative-part-program} first gives
$N_\sigma(D_V)=F_V$. Lemma~\ref{prog:negative-part-pullback} then gives
$N_\sigma(\pi^*D_V)=\pi^*F_V$. The new total space is used only as a
carrier for this equality and for pullbacks of positive currents;
no effective boundary on it is needed. Rename these spaces $V$
and $S$. No additional relative program is needed for this modification.

The vertical-divisor descent and negative-part lemmas above now
give actual effective divisors $F_S,F_V$ with
\[
 F_V=g^*F_S,\qquad
 D_S=(a+1)\alpha_S+[F_S],\qquad N_\sigma(D_S)=F_S.
\]
Apply $\mathsf M_{d-1}$ to the generalized pair with adjoint $D_S$.
On a common resolution $p:W\to S$, $q:W\to S_m$ of the resulting
good model, write
\[
 p^*D_S=q^*D_m+[E],\qquad E\geq0
 \text{ exceptional over }S_m.
\]
The negative-part comparison proved above gives
$E=p^*F_S$, and hence
\[
 q^*D_m=(a+1)p^*\alpha_S.
\]
Since $D_m$ is semiample, $\alpha_S$ is semiample after descent
through $p$. Pulling back by $g$ and using the
$\alpha$-trivial initial program gives semiampleness of
$\alpha$ on $X$. The target is compact K\"ahler, and the
resulting map is Moishezon, as follows by taking common projective
modifications of the maps just constructed. The comparison uses only the chosen lower good model; no
lower-dimensional program is lifted to $V$.

\end{proof}

\begin{proposition}[A negative ray with an actual detector]
\label{prog:ordinary-step}
Assume $\mathsf B_d$. Let $A$ be a gklt adjoint in dimension at most
$d$ on a globally strongly $\mathbb Q$-factorial compact K\"ahler
space, with globally nef carrier data. Let $R$ be an $A$-negative
extremal ray of the full analytic cone. If an actual global rational
line $L$ has $L\cdot R<0$, the ray has a projective contraction;
a small contraction has its detected flip. The working models retain
the strong factoriality and compact K\"ahler properties.
In particular this applies to an ordinary rational adjoint
$A=c_1(K_X+B)$ with detector $K_X+B$, and to ordinary real boundaries
by Proposition~\ref{prog:ordinary-replacement}.
\end{proposition}
\begin{proof}
The local polyhedrality of the negative cone makes $R$ an exposed ray.
Choose a nef support $\lambda$ whose null analytic face is $R$.
Normalize the analytic cone by a K\"ahler class. On a neighbourhood
of the point representing $R$, the class $-A$ is positive. On the
remaining compact part of this slice, $\lambda$ has a positive
minimum. Consequently, for a sufficiently large $b$, the class
$\omega_R=b\lambda-A$ is K\"ahler. Thus
$A+\omega_R=b\lambda$ is nef with precisely the null face $R$.
The gklt presentation with nef carrier datum increased by the pullback
of $\omega_R$ has modified-big total boundary. Apply $\mathsf B_d$
to construct its contraction. Lemma~\ref{prog:detected-projectivity}
makes the same global line $-L$ relatively ample, and
Proposition~\ref{prog:detected-flip} constructs the flip when needed.
For a divisorial contraction the exceptional-prime negativity and
integral line descent give strong factoriality on the target: pull
back a rank-one reflexive sheaf, remove its degree by a rational
multiple of the exceptional divisor, descend the resulting line,
and compare on the common big open. Local ordinary log-Fano
replacement gives rational singularities by relative vanishing.
The supporting target and projective flip are compact K\"ahler.
\end{proof}

\subsubsection{Relative good models and finite geography}

\begin{lemma}[Removing the preparation error]
\label{prog:remove-preparation}
Let $\mu:W\to X$ be a projective smooth preparation over $S$ with
$A_W=\mu^*A_X+[F]$, where $F\geq0$ is $\mu$-exceptional and has
positive coefficient on every added prime. If $W\dashrightarrow Y$
is a chosen good log terminal model of $A_W$ over $S$, then it induces
a good log terminal model of $A_X$ over $S$.
\end{lemma}
\begin{proof}
On a common projective resolution $p:V\to W$, $q:V\to Y$, write
$p^*A_W=q^*A_Y+[E]$ with $E\geq0$ exceptional over $Y$.
Put $r=\mu p$ and $D=E-p^*F$. Then
\[
 [D]=r^*A_X-q^*A_Y,\qquad r_*D=r_*E\geq0.
\]
The divisor $-D$ is $r$-nef because $A_Y$ is nef over $S$.
Negativity gives $D\geq0$, hence $E\geq p^*F$. Every added prime
is therefore exceptional over $Y$. This proves nonextraction from
$X$ and gives its effective exceptional comparison with $A_Y$.
Strictness at any contracted prime of $X$ follows from strictness
for the chosen log terminal model of $W$, since $F$ has coefficient
zero there. The same semiample endpoint makes the model good.
\end{proof}

\begin{lemma}[Recovering the uncontracted original primes]
\label{prog:weak-to-terminal}
Let a gklt pair on $X$ have a nonextracting weak good model
$X\dashrightarrow Y$ over $S$. Assume $Y$ is compact K\"ahler,
globally strongly $\mathbb Q$-factorial and has the exceptional
splitting of Lemma~\ref{prog:detected-cohomology}. Then a projective
crepant extraction $Y'\to Y$ gives a good log terminal model of the
original pair. The exceptional splitting persists on $Y'$.
\end{lemma}
\begin{proof}
On a common resolution write the actual comparison
$p^*A_X=q^*A_Y+[E]$, where $E\geq0$ is $q$-exceptional.
There are finitely many primes of $X$ contracted by its marking.
Let $\mathcal P$ consist of those with coefficient zero in $E$.
Take a projective log resolution $r:W\to Y$ which contains all
these valuations and dominates the fixed nef carrier. In the crepant
structure boundary $B_W$, the coefficient of each member of
$\mathcal P$ is its original effective boundary coefficient, hence
lies in $[0,1)$. Keep these coefficients unchanged. Increase each
other $r$-exceptional coefficient, if necessary from a negative
value, to a number in $(\max\{0,b_Q\},1)$. Keep all nonexceptional
coefficients unchanged. The resulting effective klt boundary
$\Gamma_W$ satisfies
\[
 [K_W+\Gamma_W]+\mathbf M_W=r^*A_Y+[F],
\]
where $F\geq0$ has precisely the unwanted exceptional primes as its
support. The inherited splitting gives actual real-line representatives
for every class modulo $Y$. Since $r$ is birational, the nef carrier
datum is relatively big; the adjoint is relatively represented by
$F$ and is pseudo-effective. Apply
Proposition~\ref{prog:projective-good-model} over $Y$.
Its projective endpoint $r':Y'\to Y$ has an effective relatively nef
exceptional trace $F'$, so negativity gives $F'=0$. A prime outside
$\operatorname{Supp}F$ cannot be contracted by an $F$-negative
birational program: on a general curve through its generic point the
effective divisor has nonnegative degree. Thus all members of
$\mathcal P$ survive, and all unwanted exceptional primes disappear.
No new prime is extracted by the program. It follows that $r'$ is
crepant and that the marking from $X$ contracts exactly primes with
strictly positive original comparison coefficient. It is nonextracting
and is therefore a log terminal model. The pullback of the semiample
adjoint on $Y$ makes it good. Finally apply
Corollary~\ref{prog:extraction-cohomology}.
\end{proof}

\begin{lemma}[A detector throughout the big-adjoint construction]
\label{prog:big-good-model}
Assume $\mathsf B_d$ and $\mathsf F_{d-1}$. The big-adjoint
construction of \cite[Section~5.1]{HX} gives a chosen good model in
dimension $d$ using only detected steps. It also gives the case of
$\mathsf M_d$ in which $A+c f^*\omega_S$ is big for some $c>0$,
with the map to $S$ retained throughout.
\end{lemma}

\begin{proof}
After the smooth preparation in \cite[Proposition~5.12]{HX}, write the
adjoint as
\[
 A=D+\omega,\qquad D\geq0,\quad \omega\text{ K\"ahler},
 \qquad \operatorname{Supp}D=N_+(A).
\]
Here $N_+(A)$ denotes the reduced divisor with support
$\bigcap_{0<\epsilon\ll1}\operatorname{Supp}N_\sigma(A-\epsilon\omega)$;
this support is constant for sufficiently small positive $\epsilon$
and is independent of the chosen Kähler class $\omega$
\cite[Lemma~5.3]{HX}. At a nonterminal scaling step,
\[
 A_i\cdot R_i<0,\qquad
 (A_i+t_i\omega_i)\cdot R_i=0,\qquad t_i>0.
\]
Consequently $\omega_i\cdot R_i>0$ and $D_i\cdot R_i<0$.
A component of the actual effective divisor $D_i$ supplies a global
rational Cartier detector. The supporting contraction follows from
$\mathsf B_d$, and its projectivity and flip follow from
Proposition~\ref{prog:detected-flip}. This argument does not require
$\omega_i$ to remain nef.

To apply $\mathsf B_d$ to a gdlt presentation, lower its finitely many
floor coefficients on the initial smooth carrier and add the same small
class to the K\"ahler nef datum. The adjoint does not change, and the
resulting gklt presentation persists through the same negative steps.
The original gdlt presentation is retained for adjunction to the floor.
Thus the special-termination proof of \cite[Lemma~5.8]{HX} applies with
$\mathsf F_{d-1}$ on its smooth lower-dimensional carrier. Its other
inputs are adjunction, discrepancy monotonicity, and the local
Cartier-index calculation. The last neighbourhood-isomorphism step
uses the evaluation-ideal support argument in
the \hyperref[prog:special-support]{support comparison above}: the comparison of
the first detected flip has support equal to the full inverse image
of its exceptional locus, and later comparisons dominate it.
One does not infer disjointness of images merely from a vanishing
restriction of an arbitrary exceptional divisor. The theta induction of
\cite[Proposition~5.11]{HX} now gives a weak model. Applying
$\mathsf B_d$ to its nef adjoint makes it good, and
Lemma~\ref{prog:weak-to-terminal} gives the log terminal model.
All its steps are detected or belong to the already-projective
relative extraction in that lemma.

For completeness, the negative-part identity needed in this argument
is only
\[
 N(p^*\xi+E)=N(p^*\xi)+E\quad(E\geq0\text{ $p$-exceptional}),
\]
not the stronger formula $N(p^*\xi)=p^*N(\xi)$ for an arbitrary
pseudo-effective class. The comparison with a nef endpoint identifies
its exceptional error with the negative part. The formula for
$N_+$ follows by writing a K\"ahler form upstairs as
$p^*\omega-F$, with $F$ positive on every exceptional prime, and applying
the displayed identity to
$p^*(\xi-\epsilon\omega)+E+\epsilon F$. Pushforward identifies the
nonexceptional coefficients; $\epsilon F$ supplies all exceptional
primes. These are exactly the uses in the theta argument.

Now put $P=f^*\omega_S$ and suppose $A+cP$ is big for some $c$.
Fix $\delta>0$ smaller than the $\omega_S$-degree of every compact
integral curve on $S$ and enlarge $c$ so $c>4d/\delta$.
Use the following direct preparation, which keeps the original
unshifted globally nef datum. On a smooth projective carrier
$\nu:W\to X$ write
\[
 \nu^*(A+cP)=[G]+\omega,\qquad G\geq0,\quad\omega\text{ K\"ahler},
 \qquad [K_W+B_W]+M_W=\nu^*A,
\]
with the non-K\"ahler support property used in the big-class
preparation. Here $\operatorname{Ex}(\nu)$ below denotes the reduced
divisor formed by all $\nu$-exceptional primes. Write
$B_W=B_W^+-B_W^-$. For small $\epsilon>0$ put
\[
 \begin{split}
 B'&=B_W^++\epsilon B_W^-+
                   \epsilon(G+\operatorname{Ex}(\nu)),\\
 M'_W&=M_W+c\nu^*P+\epsilon\omega,\\
 b&=B_W^++\frac{\epsilon}{1+\epsilon}\operatorname{Ex}(\nu).
 \end{split}
\]
The new boundary is klt and $M'_W$ is K\"ahler. Direct addition gives
\[
 A'=[K_W+B']+M'_W
   =(1+\epsilon)\bigl([K_W+b]+M_W+c\nu^*P\bigr).
\]
Moreover $B'\geq b$. In every theta branch the cumulative added
boundary satisfies $0\leq C\leq1-B'$, so
\[
 b+\frac{C}{1+\epsilon}
 \leq b+\frac{1-b}{1+\epsilon}<1.
\]
Thus the unshifted presentation
$[K_W+b+C/(1+\epsilon)]+M_W$ is genuinely gklt with the
\emph{original} globally nef datum. This is the presentation used
for the horizontal length bound. No base form is subtracted from
an arbitrary newly prepared nef datum.

A horizontal negative ray would have $P$-degree at least $\delta$;
the bound $0<-A_{\rm unshifted}\cdot C\leq4d$ would contradict
$c\delta>4d$. Every constructed step is therefore over $S$.
Its projective connected fibres are contracted by $f$, so $f$ factors
through its base and through the flip. The final finite extraction is
also over $S$.

The shifted nef endpoint is semiample by $\mathsf B_d$. Its zero face
contains no horizontal curve: decompose such a curve by the unshifted
adjoint's cone theorem; nefness of the shifted class forces all
summands into the zero face, and a summand of positive $P$-degree again
contradicts the length bound. The shifted class is big, so the
connected semiampleness map is bimeromorphic. Its connected fibres are Moishezon and are connected
by chains of compact curves. The preceding zero-face argument makes
the map to $S$ constant on each such curve, hence on every fibre.
The analytic factorization lemma therefore factors the map to $S$
through the semiampleness map. Subtracting $c\omega_S$ from the
normalized class on its target gives the required relative K\"ahler class for the unshifted adjoint. The effective
resolution error from the preparation is removed by the usual
negativity comparison, giving a model of the original data.

The forward traces and exceptional splittings persist by
Lemma~\ref{prog:detected-cohomology}. In the last extraction,
the newly extracted prime classes move from the old exceptional span
into the pullback span; they are actual rational Cartier classes.
Both-sign negativity proves directness as before.
\end{proof}

\begin{lemma}[Changing the base of a relatively trivial adjoint]
\label{prog:graph-good-model}
Suppose $g:X\to Y$ is projective, its general fibre is rationally
connected, and a gklt adjoint satisfies $A_X=g^*A_Y$. Assume the
prepared nef datum is K\"ahler on a smooth carrier. Let
$Y\dashrightarrow Y_m$ be the chosen nonextracting good model
supplied by $\mathsf M_{d-1}$ over $S$, with its exceptional
cohomology splitting, and let $\dim Y<d$. Then a single projective relative construction gives a weak good model of $(X,A_X)$ over $S$ with that splitting.
Lemma~\ref{prog:weak-to-terminal} makes it log terminal. No step
of the base program is lifted.
\end{lemma}

\begin{proof}
Take a common projective smooth resolution $T\to Y$, $T\to Y_m$
and resolve the main component of $X\times_Y T$, further dominating
the prepared nef carrier; denote the resulting smooth space by $V$.
Its morphism $h:V\to Y_m$ is projective. Before constructing a
program, verify algebraicity for all classes. The projective RC map
$V\to T$ and Lemma~\ref{prog:relative-hodge} express every class
modulo $T$ by actual real lines. The splitting of the chosen $Y_m$
expresses classes from $T$, modulo $Y_m$, by exceptional divisor
classes. Hence
\[
 H^{1,1}_{\rm BC}(V)=c_1(\operatorname{Pic}(V)\otimes\mathbb R)
                         +h^*H^{1,1}_{\rm BC}(Y_m).
\]
Lemma~\ref{prog:projective-cone} identifies the relative cone
with the full analytic face. These properties and the exceptional
splitting persist along the projective relative steps by
Lemma~\ref{prog:detected-cohomology}.

Choose the effective gklt preparation on this final carrier. Its
adjoint has the identity
\[
 A_V=h^*A_{Y_m}+[G],\qquad G=G_{\rm base}+F\geq0,
\]
where $G_{\rm base}$ is pulled back from the actual effective base
comparison and $F$ is exceptional over $X$, with positive coefficient
on every added prime. Its nef datum is nef and big on the prepared
carrier. Modulo $Y_m$ it is the actual real-line class
$[G]-c_1(K_V)-[B_V]$. Thus every hypothesis of
Proposition~\ref{prog:projective-good-model} holds, with the effective
relative adjoint $G$. Denote its global relative endpoint by $V_m$
and the strict transform of $G$ by $G_m$.

On a common resolution $p:W\to V$, $q:W\to V_m$, the actual adjoint
comparison has effective $q$-exceptional error $E$. The difference
$p^*G-q^*G_m-E$ has zero Bott--Chern class and is $q$-exceptional,
because $q_*p^*G=G_m$. Negativity in both signs makes it zero:
\[
 p^*G=q^*G_m+E.                                      \tag{*}
\]
Choose a big common-isomorphism open $U\subset Y_m$ for
$Y\dashrightarrow Y_m$. Such an open exists because that map is
nonextracting. On $U$ we have $G_{\rm base}=0$. Let
$r:W_U\to X_U$ be the induced projective birational map. The divisor
$q^*G_m$ is $r$-nef, since the endpoint adjoint is nef over $U$ and
the base class has zero degree on $r$-fibres. Equation $(*)$ gives
\[
 r_*q^*G_m=-r_*E\leq0,
\]
since $F$ is exceptional over $X$. Negativity yields $q^*G_m\leq0$.
Effectivity yields $G_m|_U=0$. This step removes the possible
\emph{horizontal} components of $F$; they were not assumed very
exceptional at the start.

Now $G_m$ is supported over codimension at least two in $Y_m$.
It is nef over $Y_m$, so very-exceptional negativity gives $G_m=0$.
The endpoint adjoint is the pullback of $A_{Y_m}$, hence is nef and
good over $S$. Its relative canonical morphism is Moishezon: it is
the connected factor of the composite of the projective morphism to
$Y_m$ and the lower-dimensional Moishezon canonical morphism.
Lemma~\ref{prog:remove-preparation} removes the added
source-exceptional errors and proves nonextraction. All these operations are over $S$.

The all-class invariant was established before the relative program;
its preservation follows at each step from the same cone identification
and forward cohomology splitting. If the construction has contracted
any original prime crepantly, apply
Lemma~\ref{prog:weak-to-terminal} to retain it.
\end{proof}

\begin{lemma}[The global rational quotient over a fixed base]
\label{prog:relative-model}
Assume $\mathsf M_{d-1}$ and the case of $\mathsf M_d$ in which a
base shift makes the adjoint big. Then $\mathsf M_d$ holds when no
class $A+c f^*\omega_S$, $c\geq0$, is big.
\end{lemma}

\begin{proof}
First take the usual smooth carrier with effective gklt boundary
and adjoint $\mu^*A+[F_0]$, where $F_0$ is effective exceptional
with full exceptional support. Its boundary-plus-nef trace is big and its adjoint
is not big, so its canonical class is not pseudo-effective. Take
its global MRC and resolve the MRC graph and the fixed carrier,
obtaining a projective map to a smooth compact K\"ahler $Z$ with
$K_Z$ pseudo-effective.
Perform the K\"ahler-nef-part preparation \emph{after} these graph
resolutions, preserving the maps to $Z$ and $S$. This order ensures
that the final scaling datum is K\"ahler, rather than merely the nef
pullback of an earlier K\"ahler datum. Write the resulting adjoint as
\[
 A_0=K_U+B_U+M_U=\mu^*A+[F],
 \qquad F\geq0\text{ exceptional},\qquad M_U\text{ K\"ahler}.
\]
The global modified-big hypothesis makes a preliminary base shift
unnecessary. Relative pseudo-effectivity is preserved.
No base shift of $A_0$ is big, since pushing a big such class down
would make a base shift of $A$ big. The fixed morphism $h:U\to Z$
is projective and has RC general fibres. All these preparations
precede the next, fresh base shift.

Put $P=f_U^*\omega_S$ and choose $c>2d/\delta$, where every
integral curve on $S$ has $\omega_S$-degree greater than
$\delta>0$. Set
\[
 A_c=A_0+cP=D+N,\qquad D=K_U+B_U,\quad N=M_U+cP.
\]
Since $A_c$ is not big globally, relative-canonical positivity
\cite[Theorem~2.48(2)]{HX}, together with pseudo-effectivity of $K_Z$,
shows that $A_c$ is not big over $Z$, whether or not it is globally
pseudo-effective. The class $N$ is big over $Z$. Consequently $D$
cannot be pseudo-effective over $Z$: otherwise $D+N=A_c$ would be
big over $Z$. The smooth projective RC Hodge projection represents
every class modulo $Z$ by an actual real line. The ordinary global klt adjoint $D$ is therefore relatively
non-pseudo-effective. The projective ordinary Mori construction of
Corollary~\ref{prog:projective-mori}, with scaling of the actual
relatively ample real line represented by $N$, terminates in a Mori
fibre space. Lemma~\ref{prog:relative-hodge} and
Lemma~\ref{prog:projective-cone} identify every relative ray
with a full analytic ray. Its nef b-data retain the fixed globally
nef carrier $U$.

First suppose $A_c$ is not pseudo-effective globally. The retained
relative-canonical positivity theorem on $h$, together with
pseudo-effectivity of $K_Z$, then says that $A_c$ is not
pseudo-effective over $Z$. The terminating relative program ends
in a Mori contraction at a threshold $\tau>1$. At every step its
scaled zero ray has $A_c$-degree negative. Inductively the map to
$S$ is retained: if such a ray were horizontal, its $P$-degree
would exceed $\delta$, while the gklt cone theorem for $A_0$ gives
a rational generator of $-A_0$-degree at most $2d$, contradicting
$A_c\cdot R\leq0$. Once $P\cdot R=0$, the step is negative or
crepant for $A_0$, so its gklt presentation persists. The same
argument applies to the final Mori ray. This would be an
$A_c$-negative Mori contraction over $S$, contradicting the
preserved pseudo-effectivity over $S$. Hence $A_c$ is
pseudo-effective globally.

It is not big globally. Relative-canonical positivity now implies
that it is not big over $Z$. Its pseudo-effective threshold for
$D+tN$ over $Z$ is exactly one: for $t>1$ the class is big, while
pseudo-effectivity for $t<1$ would make $A_c$ big because $N$ is
big over $Z$. The same finite relative program therefore ends in
a Mori contraction $g:U'\to Y$ at threshold one.
All birational steps, including any steps at threshold one, and
the final zero ray are over $S$ by the preceding length argument
with $A_c\cdot R\leq0$. For a zero ray with positive $P$-degree,
$A_0$ would again have degree less than $-c\delta$.
Thus $Y$ carries a map to $S$ and $\dim Y<d$.
The class $A_c$ has zero degree on every $g$-contracted curve.
This is an ordinary projective log-Fano contraction: rationalize its
effective real ordinary boundary, preserving klt and relative
antiampleness. Relative vanishing gives rational singularities on
$Y$, which is compact K\"ahler because it is projective over $Z$.
Lemma~\ref{prog:bc-descent} therefore gives an actual equality
$A_c=g^*A_Y$ in Bott--Chern cohomology. This descent precedes the
base presentation. Apply Proposition~\ref{prog:base-presentation}
to $g$ and the whole adjoint $A_c$. The negative or crepant steps
preserve its effective gklt data. Its nef part is the trace of
$M_U+cP$ from the fixed carrier; on a common carrier it is nef and
big, so its total boundary is globally modified big. The proposition
gives effective gklt data for $A_Y$ with globally nef carrier and
globally modified-big total boundary. Relative pseudo-effectivity
descends through the projective surjection $g$, so $\mathsf M_{d-1}$
applies over $S$, with the base shift retained. For the source data in
Lemma~\ref{prog:graph-good-model}, use
Lemma~\ref{prog:base-presentation-preparation} to choose a K\"ahler
nef datum on a smooth carrier while keeping $A_c$ fixed. Apply that
graph lemma to the chosen lower-dimensional output.

Adjoint-preserving preparation is compatible with the original pair
in this comparison. On a common resolution of a nonextracting good
output, keep the original nef datum and subtract the effective model
comparison divisor from the original structure divisor. The resulting
structure divisor has the same adjoint as the output; its trace is
effective by nonextraction, and its discrepancies have not decreased.
Strictness on contracted original primes is preserved. Thus the
prepared presentation supplies a comparison for the original data as
well.

Steps at threshold one can be $A_c$-crepant divisorial contractions.
Consequently the composite constructed so far need only be a weak
good model of the original prepared adjoint. Apply
Lemma~\ref{prog:weak-to-terminal} to extract the finitely many original
primes with zero comparison coefficient. This gives the strict log
terminal model and preserves the cohomology invariant. The shift
changes neither relative comparisons nor relative nefness.
Lemma~\ref{prog:remove-preparation} removes the full-support
preparation error and proves the claimed $\mathsf M_d$.
\end{proof}

\subsubsection{Finite geography and closing the induction}

\begin{lemma}[Canonical models first, weak models second]
\label{prog:finite-geography}
Assume $\mathsf B_d$ and $\mathsf M_d$. Then $\mathsf F_d$ holds for compact polytopes
of globally nef b-data on a fixed carrier and globally modified-big
boundary, over the fixed base.
\end{lemma}

\begin{proof}
First prove, by induction on the parameter-polytope dimension,
polyhedrality of the effective locus, finite canonical-model chambers,
and finitely many \emph{chosen} terminal models on those chambers,
following \cite[Theorem~5.15]{HX}.
The zero-dimensional case uses $\mathsf M_d$ and uniqueness of the
canonical model; it does not yet assert finiteness of all weak models.

At a central point choose $X\dashrightarrow X_m\to X_c$ over
$S$. The transport invariant in $\mathsf M_d$ supplies the traces
of the whole nearby polytope on $X_m$. Shrink so the finitely many
strict exceptional inequalities remain strict. Add a common pullback
from $S$ so that the central canonical class on $X_c$ is K\"ahler
absolutely. On the boundary of the parameter polytope, apply the
smaller-parameter induction \emph{over $X_c$}; its existence input
is $\mathsf M_d(X_c)$, just proved independently. No projectivity of
$X_m\to X_c$ or Hodge projection over the possibly singular $X_c$
is used.

The central adjoint is zero over $X_c$, so the relative effective
locus is the radial cone on its boundary effective locus. Apply the
cone-length estimate of \cite[Claim~5.1]{HX} on the finitely chosen
\emph{gklt terminal models} $h_i:Y_i\to X_c$, all of dimension $d$.
Write $A_{u,i}$ for the boundary-parameter adjoint, nef over $X_c$,
and $\psi_i:Y_i\to Z_i$ for its relative canonical morphism.
The cone theorem is not applied to the possibly lower-dimensional
canonical targets $Z_i$.

Let $H_i=h_i^*\omega_c$ and choose $\delta>0$ below the
$\omega_c$-degree of every compact integral curve in $X_c$. Fix $\bar\lambda<\delta/(\delta+2d)$ and for
$0<\lambda\leq\bar\lambda$ put
\[
 A_{\lambda,i}=(1-\lambda)H_i+\lambda A_{u,i}.
\]
The same cone estimate shows more than nefness: for a sufficiently
small $\epsilon>0$, uniform on the finite list and the chosen radial
interval, $A_{\lambda,i}-\epsilon H_i$ is nef. Indeed on a
horizontal $A_{u,i}$-negative rational ray its degree is greater than
$(1-\lambda-\epsilon)\delta-2d\lambda>0$; on the nonnegative
part of the cone and on the vertical cone it is nonnegative.
It follows for every class $\xi\in\overline{\mathrm{NA}}(Y_i)$ that
\[
 A_{\lambda,i}\cdot\xi=0
 \quad\Longleftrightarrow\quad
 H_i\cdot\xi=A_{u,i}\cdot\xi=0.
\]
In particular its null curves are exactly those contracted by
$\psi_i$. The central data are crepant through these maps over $X_c$;
thus the convex data defining $A_{\lambda,i}$ remain gklt with
globally modified-big boundary. Apply $\mathsf B_d$ to $A_{\lambda,i}$. Its connected
semiampleness map and $\psi_i$ both have Moishezon fibres, and each
connected such fibre is connected by chains of compact curves.
Equality of their contracted curves makes each map constant on the
fibres of the other. The analytic factorization lemma identifies
the two maps. Thus the selected relative canonical model is also
the canonical model over $S$ on this radial interval.
This proves local polyhedrality and the finite local list; compactness
gives the global finite list of canonical models and chosen terminal
models.

\paragraph{Uniform preparation preserving every weak model.}
\label{prog:uniform-preparation}
Write $A_u=[K_X+B_u+\mathbf M_{u,X}]$ for the adjoints in the
parameter polytope. We first record why a weak model of these gklt data
extracts no divisor. On a common resolution $p:U\to X$, $q:U\to Y$
put $D=p^*A_u-q^*A_{u,Y}$, using the actual structure-boundary
comparison. The weak-model discrepancy inequalities give $p_*D\geq0$,
and $-D$ is $p$-nef because $A_{u,Y}$ is nef over the fixed base.
Negativity therefore gives $D\geq0$. An extracted prime on $Y$ has
coefficient one in the weak-model boundary, whereas its coefficient in
the structure boundary of the gklt source is less than one; it would
give a negative coefficient of $D$. Hence there is no such prime.
Consequently $D$ is $q$-exceptional and the target is itself gklt.

Fix a parameter $u_0$. The boundary plus nef part is globally
modified big by hypothesis. The modified-big preparation
\cite[Lemma~2.23]{HX} gives, on a fixed projective smooth carrier
$a:W\to X$, an effective klt boundary $\Gamma$, a K\"ahler class
$\gamma_0$, and an effective $a$-exceptional real divisor $F$ such that
\begin{equation}\label{prog:eq:uniform-preparation}
             [K_W+\Gamma]+\gamma_0=a^*A_{u_0}+[F].
\end{equation}
Here is the effectivity detail in this use of the preparation. Its new
nef datum dominates a K\"ahler form on a carrier. After resolving the
new structure boundary, choose an effective exceptional divisor whose
negative is relatively ample and subtract a sufficiently small multiple
from the pulled-back datum to make it K\"ahler. Add the same multiple
to the structure boundary,
then replace its negative coefficients by zero and, if needed, add
small positive coefficients on the remaining exceptional primes.
All coefficients stay below one. Only exceptional coefficients have
been increased; their increase is the divisor $F\geq0$ in
\eqref{prog:eq:uniform-preparation}.

On a sufficiently small closed polyhedral neighbourhood of $u_0$ put
\[
                \gamma_u=\gamma_0+a^*(A_u-A_{u_0}).
\]
These classes remain K\"ahler by openness, and the same $\Gamma$ and
the same $F$ satisfy
\[
                  [K_W+\Gamma]+\gamma_u=a^*A_u+[F].
\]
Thus this is one affine family of gklt pairs on one smooth carrier,
with a fixed effective boundary and K\"ahler nef data.

Let $\psi:X\dashrightarrow Y$ be \emph{any} weak model at a parameter
$u$ in this neighbourhood, and resolve the induced map $W\dashrightarrow
Y$ by $r:T\to W$, $q:T\to Y$. Its original comparison is
\[
              (ar)^*A_u=q^*A_{u,Y}+E_u,
              \qquad E_u\geq0\text{ is }q\text{-exceptional}.
\]
Since $\psi$ is nonextracting, every $a$-exceptional prime is exceptional
over $Y$. Hence $r^*F$ is also $q$-exceptional, and
\begin{equation}\label{prog:eq:prepared-weak}
 r^*\bigl([K_W+\Gamma]+\gamma_u\bigr)
       =q^*A_{u,Y}+E_u+r^*F.
\end{equation}
This proves that $W\dashrightarrow Y$ is a weak model of the prepared
pair. In detail, subtract the effective exceptional error on the right
from its structure boundary on $T$ and push down to $Y$. The resulting
boundary is precisely the transform of $\Gamma$, it is effective, the
adjoint is the already existing Bott--Chern class $A_{u,Y}$, and all
discrepancies are at least those of the prepared gklt pair. This uses generalized data at the level of currents. On $T$ keep
the nef datum $r^*\gamma_u$, and subtract the displayed effective
$q$-exceptional error from the prepared structure boundary. Pushing
forward by $q$ gives the boundary $\Gamma_Y$ and the trace of this
fixed nef b-class. Their sum with the canonical current represents
the already existing locally exact adjoint $A_{u,Y}$; the resolution
identity defines its structure boundary. The generalized discrepancies
have not decreased. This argument neither requires $K_Y+\Gamma_Y$ to
be real Cartier nor asserts that the nef trace is separately a
Bott--Chern class. It therefore applies to every K\"ahler weak target,
including the adjunction strata used in special termination.

\paragraph{Every weak model occurs in the finite canonical list.}
Choose a small box in $H^{1,1}(W,\mathbb R)$ which spans that vector
space and whose addition to the compact family $\{\gamma_u\}$ stays
inside the K\"ahler cone. For any $Y$ above choose a K\"ahler class
$\eta_Y$ and define $\eta_W=r_*q^*\eta_Y$. Smooth-source cohomology
puts $\eta_W$ in $H^{1,1}(W,\mathbb R)$ and gives an actual
$r$-exceptional real divisor $D_\eta$ with
\[
                    r^*\eta_W-q^*\eta_Y=[D_\eta].
\]
Its negative is $r$-nef, so $D_\eta\geq0$ by exceptional negativity.
As $W\dashrightarrow Y$ extracts no divisors, $D_\eta$ is also
$q$-exceptional. Scale $\eta_Y$ by a sufficiently small positive number
so that $\eta_W$ lies in the fixed box. Adding this identity to
\eqref{prog:eq:prepared-weak} yields an effective
$q$-exceptional comparison with target class $A_{u,Y}+\eta_Y$, which
is K\"ahler over the original base. Thus $W\dashrightarrow Y$ is the
canonical model of a member of this one enlarged prepared polytope.
Its canonical-model list is finite by the first part of the geography
proof. The induced list of maps from $X$ is therefore finite as well.
A finite cover of the original compact parameter polytope completes
the proof for all weak models.

\end{proof}


\begin{theorem}[The restricted package in every dimension]
\label{prog:restricted-package}
The assertions $\mathsf B_d$, $\mathsf C_d$, $\mathsf M_d$ and
$\mathsf F_d$ hold in every finite dimension, for the globally nef
carrier data and globally modified-big boundary traces specified above.
All birational steps used to prove them are detected steps or belong
to an already-projective relative construction with the stated
relative-line algebraicity. A general semiample morphism is asserted
to be Moishezon; projectivity is asserted only when the construction
provides an actual relatively ample line.
\end{theorem}
\begin{proof}
The assertions in dimensions zero and one follow from degrees and
factorization of maps of compact curves. Suppose the package holds
in dimensions less than $d$. Lemma~\ref{prog:base-point-free-induction}
first proves $\mathsf C_{d,\mathrm{big}}$ and then $\mathsf B_d$.
The former uses $\mathsf C_{d-1}$; the nonbig part of the latter
uses only a chosen $\mathsf M_{d-1}$ model and the explicit
negative-part comparison. In particular, no lower program is lifted.

With $\mathsf B_d$ available, Proposition~\ref{prog:ordinary-step}
constructs detected negative steps in dimension $d$.
Lemma~\ref{prog:big-good-model}, using the already known
$\mathsf F_{d-1}$, gives the big-shift case of $\mathsf M_d$.
Lemma~\ref{prog:relative-model} gives its remaining case, using
$\mathsf M_{d-1}$ and the projective graph construction. This proves
$\mathsf M_d$ over every proper compact K\"ahler base, with its
chosen-model invariant. Lemma~\ref{prog:finite-geography}, whose
hypotheses are \emph{both} $\mathsf B_d$ and $\mathsf M_d$, now
proves $\mathsf F_d$.

It remains to prove $\mathsf C_d$ when its nef NQC adjoint $\alpha$
is not big. First use
Lemma~\ref{prog:base-presentation-preparation}, preserving the
adjoint $\alpha$ and making the nef datum K\"ahler on a smooth
carrier. Its multiplier-ideal inclusion allows the given contraction
to restrict to the prepared non-klt closed subspace by
Lemma~\ref{prog:base-emc}. Next take the projective dlt preparation
in \cite[Section~6]{HX}, keeping its crepant structure data. The
working space is klt and globally strongly $\mathbb Q$-factorial.
The retained crepant structure boundary need not be lc; it is
effective, including the extracted primes of coefficient at least one. On a common carrier the nef datum is
nef and big. Consequently the total boundary trace is modified big,
and $\alpha-c_1(K_X)$ is a big Bott--Chern class. The contraction on
the non-klt closed subspace transports through this preparation by
Lemma~\ref{prog:multiplier-transport}. Choose an NQC expression with
positive coefficients and integral multiples of its rational
degree-two summands. Corollary~\ref{prog:uniform-canonical-parameter}
provides a single number $t$ such that each negative ray of
$K_X+t\alpha$ is $\alpha$-trivial. Choose a K\"ahler scaling class
$\omega$ with $K_X+t\alpha+\omega$ nef. Every selected negative
ray is ordinary $K_X$-negative; the actual rational canonical line
is its detector. Proposition~\ref{prog:ordinary-step} supplies its
projective contraction or flip, and the fixed integral grid persists.
The descended $\alpha$ remains nef by projective descent of nefness.

The driving class is non-pseudo-effective throughout, by the
corollary. Its initial pseudo-effective scaling threshold is positive;
choose a cutoff $\varepsilon>0$ below that threshold. On the fixed
initial carrier the compact segment of globally nef data
\[
                    t\alpha+[\varepsilon,1]\omega
\]
has modified-big total boundary. At its scaling parameter each working
model is a marked weak model of a member of this segment. Each selected
step is negative for the fixed driving adjoint, so its discrepancies
do not decrease and increase at a valuation affected by that step.
An infinite sequence would repeat a marked model by $\mathsf F_d$;
the cumulative nonzero effective comparison for the fixed adjoint
would then contradict equality of that marking. Hence this program
terminates. Non-pseudo-effectivity excludes a nef endpoint, leaving
an $\alpha$-trivial ordinary Mori contraction $f:X'\to Z$.

The structure data of the prepared contraction problem remain
crepant throughout this $\alpha$-trivial program. Their effective
traces and globally nef and big carrier datum therefore persist.
Lemma~\ref{prog:multiplier-transport} transports the contraction on
the non-klt closed subspace through each birational step.
Lemma~\ref{prog:bc-descent} gives
\[
                         \alpha_{X'}=f^*\gamma
\]
in Bott--Chern cohomology. The class $\gamma$ is nef by projective
descent; Lemma~\ref{prog:integral-h2} descends the fixed rational
NQC summands, and projective multisections preserve their
nonnegative degrees.

Apply Corollary~\ref{prog:base-emc-induction} to this projective
Mori contraction. In its notation the prepared non-klt subspace
$N_1$ has Stein space $V(f_*J_1)$ over $Z$. If $N_1$ dominates
$Z$, this Stein space is $Z$ itself and its inherited contraction
already supplies the contraction of $\gamma$. Otherwise the base
presentation has multiplier ideal $J_Z\supset f_*J_1$, so the
contraction restricts to $V(J_Z)$ as a closed subspace of
$V(f_*J_1)$. Since $\dim Z<d$, the assertion $\mathsf C_{d-1}$
then gives the contraction of $\gamma$. Pullback supplies the
contraction on $X'$. Its pullback to projective common resolutions
factors through the original space by Lemma~\ref{prog:base-emc},
because the adjoint classes are crepant. This gives the required
Moishezon contraction on the original space.
This proves $\mathsf C_d$ and completes the induction.
\end{proof}

\begin{corollary}[Finite positive parts of detected scalings]
\label{prog:detected-positive-truncation}
Let a gklt adjoint with globally nef carrier data be run with scaling
of the trace of a K\"ahler class on a fixed smooth carrier. Suppose
that every chosen negative analytic ray has an actual global rational
line detector and that the total boundary trace on each positive
scaling segment is globally modified big. Every segment whose scaling
parameters lie in a fixed interval $[\varepsilon,T]$, with
$\varepsilon>0$, is finite. If the driving adjoint is not
pseudo-effective, the scaling terminates with a Mori fibre space.
\end{corollary}
\begin{proof}
Proposition~\ref{prog:ordinary-step} constructs each step and
Lemma~\ref{prog:detected-cohomology} supplies the forward traces from
the fixed carrier. Each reached model at a scaling parameter is a
marked weak model of that parameter's adjoint. The compact parameter
segment satisfies $\mathsf F_d$, so there are finitely many such
markings. Strict increase for the fixed driving adjoint rules out
repetition. If that adjoint is not pseudo-effective, the positive
pseudo-effective scaling threshold gives a fixed positive cutoff
containing every parameter of a continuing program. The finite
program cannot end nef, hence ends with the claimed Mori map.
\end{proof}
